\documentclass[10pt,a4paper]{article}
\usepackage{amsmath,amssymb,amsthm,graphicx}
\usepackage{amsmath,amssymb,amsthm}
\newtheorem{nttheorem}{Theorem}
\newtheorem{ntdefinition}[nttheorem]{Definition}
\newtheorem{ntlemma}[nttheorem]{Lemma}
\newtheorem{ntproposition}[nttheorem]{Proposition}
\newtheorem{ntclaim}[nttheorem]{Claim}
\title{Rate amplification for non-adaptive linear locally correctable codes via dual smooth recovery}
\author{Gil Cohen, Tal Yankovitz}
\date{doi:10.4230/LIPIcs.CCC.2021.1; CC BY4.0}
\begin{document}
\maketitle
\noindent\textit{Editorial note.} Gil Cohen and Tal Yankovitz. CCC2021 article1. DOI10.4230/LIPIcs.CCC.2021.1. License CC BY4.0.\par\medskip
\noindent\textit{Editorial note.} Faithful editable transcription from the official CCC2021 CC BY4.0 published version, article1. One rate-amplification mechanism: introductory Theorem3 and formal Theorem46/Proposition55 with all of Sections4/4.1/4.2/4.3/4.4 author constructions and proofs, plus the full necessary Section5.2 axis-evasive partition construction and Claims57--63/Proposition60/Theorem62 proof. Definitions1--2/non-adaptive/linear code context and Claim82 retained. The unused probabilistic alternative Section5.1 is outside this item, as the author explicitly says it is not used; no separately counted support lemma. All43 original bibliography entries retained. Author wording/printed indices and equations are preserved without mathematical correction, including Definition25 distinct/not-necessarily-distinct wording, Claim23 z direction, Corollary49 final displayed little-o sign, Proposition60 final denominator c, Claim59 property numbering and final Theorem62 printed I indices. Author omitted details remain omitted rather than supplied: Fact24 called easy to verify, Claim82 induction, and Proposition55 final error calculation left to reader. These are the authors actual proofs, not extraction omissions. Formal54 parity wording retained as printed. All field/linearity/nonadaptivity/rate/error/smoothness/alphabet and prime-residue conditions retained. Root handles compilation/database/publication.\par\medskip
\noindent\textit{Allowed standard prerequisites (editorial).}\begin{itemize}
\item Siegel-Walfisz theorem. Original Theorem48[34,37] statement and full author Corollary49 deduction retained.
\item Basic finite-field and cyclic-group facts. Full Section5.2 author use and supporting calculations retained; no recursion into standard algebra.
\end{itemize}\medskip
\setcounter{nttheorem}{45}
\begin{nttheorem}[Author Theorem46]
There exist universal constants $m_0,c'\geq1$ such that the following holds. Let $P$ be a set of size $m\geq m_0$. Let $\mathbb F$ be a field, and let $(D^P_{\mathrm{in}},L^P_{\mathrm{in}})$ be a $(q_{\mathrm{in}},\tau_{\mathrm{in}},\varepsilon_{\mathrm{in}},\rho_{\mathrm{in}})$-dual SLR over $\mathbb F^P$. Let $0<\alpha<1$ be such that
\begin{equation}\rho_{\mathrm{in}}\geq\frac{c'}{\sqrt{\alpha\cdot\log m}}\log\left(\frac1\alpha\right).\tag{4.11}\end{equation}
Then, there exists a $(q_{\mathrm{out}},\tau_{\mathrm{out}},\varepsilon_{\mathrm{out}},\rho_{\mathrm{out}})$-dual SLR $(D^{P_{\mathrm{out}}}_{\mathrm{out}},L^{P_{\mathrm{out}}}_{\mathrm{out}})$ over $\mathbb F^{P_{\mathrm{out}}}$, with $m^\ell/2\leq|P_{\mathrm{out}}|\leq m^\ell$, where
\begin{equation}\ell=\Theta\left(\frac1{\rho_{\mathrm{in}}}\log\frac1\alpha\right),\tag{4.12}\end{equation}
having the following parameters:
\begin{align*}q_{\mathrm{out}}&\leq q_{\mathrm{in}}^{\operatorname{poly}(\ell)},\\\tau_{\mathrm{out}}&\leq q_{\mathrm{in}}^{\operatorname{poly}(\ell)}\tau_{\mathrm{in}}^\ell,\\\varepsilon_{\mathrm{out}}&\leq q_{\mathrm{in}}^{\operatorname{poly}(\ell)}(\tau_{\mathrm{in}}+\varepsilon_{\mathrm{in}}),\\\rho_{\mathrm{out}}&\geq1-\alpha.\end{align*}
\end{nttheorem}

\section*{Author code definitions and introductory main result}
\setcounter{nttheorem}{0}
\begin{ntdefinition}[Author Definition1: Locally decodable codes (LDC)]
A code $C:\Sigma^k\to\Sigma^n$ is $(q,\delta,\varepsilon)$-locally decodable if there exists a randomized algorithm $D$, called a local decoder, that is given $i\in[k]$ as input and an oracle access to $z\in\Sigma^n$, and has the following guarantee. For every $i\in[k]$, $m\in\Sigma^k$ and $z\in\Sigma^n$ such that $\operatorname{dist}(C(m),z)\leq\delta$ it holds that $\Pr[D^z(i)\ne m_i]\leq\varepsilon$. Moreover, $D$ makes at most $q$ queries to $z$.
\end{ntdefinition}
\begin{ntdefinition}[Author Definition2: Locally correctable codes (LCC)]
A code $C\subseteq\Sigma^n$ is $(q,\delta,\varepsilon)$-locally correctable if there exists a randomized algorithm $D$, called a local corrector, that is given $j\in[n]$ as input and an oracle access to $z\in\Sigma^n$, and has the following guarantee. For every $j\in[n]$, $c\in C$ and $z\in\Sigma^n$ such that $\operatorname{dist}(c,z)\leq\delta$ it holds that $\Pr[D^z(j)\ne c_j]\leq\varepsilon$. Moreover, $D$ makes at most $q$ queries to $z$.
\end{ntdefinition}
We place $z$ in the upper script in our notation $D^z(i)$ to stress that the number of symbols read from $z$ by $D$ is of importance. The parameter $q$ is called the query complexity, and $\delta$ is the local error decoding radius, in the case LDC, and local error correction radius in the case of LCC. However, we also refer to $\delta$, somewhat inaccurately, as the local distance of the code. From here on, we do not make any explicit reference to the ``global'' distance of a code and so we refer to the local distance simply as the distance. Throughout the paper, we only consider non-adaptive LDC and LCC, defined next. Informally, these are code in which the local decoder (or corrector) samples the entries to be read before the querying step takes place. Our results only hold for non-adaptive LDCs and LCCs. For ease of discussion, throughout the introduction we ignore the error parameter $\varepsilon$. More precisely, when stating our results, every LDC or LCC (both in the hypothesis as well as in the LDC or LCC guaranteed by the theorem) has constant error.
Clearly, an LCC with a systematic encoding\footnote{An encoding from messages to codewords is called systematic if the symbols of each message are embedded in its mapped codeword.} is also an LDC and so, in particular, linear LCC induce LDC.
\begin{nttheorem}[Author introductory Theorem3: Main result]
Assume one has a non-adaptive linear $(q,\delta=\Omega(1))$-LCC with block-length $n_0$ having rate $\rho=\rho(n_0)$. Then, for every integers $\ell,c\geq1$ such that $\ell^2<c<\log n_0$, one can obtain a non-adaptive linear $(q_{\mathrm{new}},\delta_{\mathrm{new}})$-LCC with block length $n\approx n_0^\ell$, having rate $\rho_{\mathrm{new}}$, where
\begin{align*}q_{\mathrm{new}}&=(cq)^{\operatorname{poly}(\ell)},\\\delta_{\mathrm{new}}&=(cq)^{-\operatorname{poly}(\ell)},\\\rho_{\mathrm{new}}&=1-(1-\rho)^\ell-O\left(\frac{\ell^2}c\right).\end{align*}
\end{nttheorem}

\section*{4 Rate amplification for dual-induced SLR}
In this section we introduce the notion of smooth locally recoverable sets (SLR) which under non-adaptive and linearity assumptions is shown to be equivalent to LCC. We consider a certain class of SLR, to which we call dual-induced SLR. These are SLR that are obtained by the dual of certain structured sets. The structure of these dual-SLR sets allows us to devise a rate amplification procedures for them. Informally, dual-SLR are sets of tuples (or linear spaces of vectors if the alphabet over which we are working is a field) in which every given entry of a tuple in the set can be recovered using only few queries and in a ``smooth'' manner, which is to say that the distribution of every query has high min-entropy.
\setcounter{nttheorem}{20}
\begin{ntdefinition}[Author Definition21: Smooth locally recoverable sets (SLR)]
Let $\Sigma,P$ be arbitrary non-empty sets. We say that $C\subseteq\Sigma^P$ is $(q,\tau,\varepsilon)$-smooth locally recoverable (SLR for short) if there exists a randomized algorithm $\operatorname{Rec}$, called a recovering procedure, that is given as input $p\in P$ and an oracle access to $c\in C$. The recovering procedure outputs either an element of $\Sigma$ or a symbol $\bot$ which is assumed not to be in $\Sigma$. The algorithm $\operatorname{Rec}$ has the following properties:
\begin{itemize}
\item For every $(c,p)\in C\times P$, $\operatorname{Rec}^c(p)$ makes at most $q$ queries to $c$.
\item For every $c\in C$ and $p,r\in P$ it holds that $\Pr[\operatorname{Rec}^c(p)\text{ queries }c_r]\leq\tau$.
\item For every $(c,p)\in C\times P$, the random variable $\operatorname{Rec}^c(p)\in\{c_p,\bot\}$, and $\Pr[\operatorname{Rec}^c(p)=\bot]\leq\varepsilon$.
\end{itemize}
We assume that for every $p\in P$ whether $\operatorname{Rec}^c(p)=\bot$ is independent of $c$, and that it is never the case that $\operatorname{Rec}^c(p)$ queries $c_p$. When $\Sigma$ is a field and $C$ is a linear subspace of $\Sigma^P$, we say that $C$ is linear. In this case, the rate of $C$ is defined as $\dim(C)/|P|$. We will mostly consider non-adaptive SLR. These are SLR in which the joint distribution of queries is independent of $c$.
\end{ntdefinition}
We remark that the notion of SLR is very similar to the notion of smooth-codes of~\cite{ref23} for LDC. We now have the following easy claim showing that SLR yield LCC and, assuming linearity, LDC.
\begin{ntclaim}[Author Claim22]
Let $C\subseteq\Sigma^P$ be a $(q,\tau,\varepsilon)$-SLR. Then, for every $\varepsilon'>0$, $C$ is a $(q,\delta,\varepsilon+\varepsilon')$-LCC with $\delta=\varepsilon'/(q\tau|P|)$. As a consequence, if $C$ is also linear then $C$ is a $(q,\delta,\varepsilon+\varepsilon')$-LDC.
\end{ntclaim}
\begin{proof}
To show that $C$ is an LCC, we devise a local corrector for $C$. Given an oracle access to $c\in\Sigma^P$, and $p\in P$ as input, the local corrector computes $z=\operatorname{Rec}^c(p)$. If $z=\bot$ then the local corrector returns some arbitrary element of $\Sigma$, and otherwise return $z$. To analyze this local corrector, let $c'\in\Sigma^P$ be such that $\operatorname{dist}(c,c')\leq\delta|P|$. Denote $B=\{p\in P\mid c_p\ne c'_p\}$. Note that conditioned on $\operatorname{Rec}^c(p)\ne\bot$, the local corrector returns $c_p$ successfully if all $q$ queries do not fall into $B$. The probability that any given query falls into $B$ is bounded above by $\tau|B|$ and so, by the union bound, the probability that some query falls into $B$ is bounded above by $\tau|B|q\leq\varepsilon'$. This proves that $C$ is a $(q,\delta,\varepsilon+\varepsilon')$-LCC. Note that linear LCC are systematic and so every linear LCC induces an LDC.
\end{proof}
In fact, for linear LCCs, the other direction also holds, meaning that such an LCC is an SLR, as we have in the following claim.
\begin{ntclaim}[Author Claim23]
Let $C\subseteq\mathbb F^P$ be a linear non-adaptive $(q,\delta,\varepsilon)$-LCC where $1-\varepsilon>1/|\mathbb F|$. Then, $C$ is a $(q,\tau,\varepsilon')$-SLR with $\tau=q/(\delta|P|)$ and $\varepsilon'=0$.
\end{ntclaim}
Before we prove the claim, we need to state the following easy to verify fact.
\begin{nttheorem}[Author Fact24]
Let $L\subseteq\mathbb F^P$ be a linear subspace, let $p\in P$, let $Q\subseteq P$ and let $x\in\mathbb F^{|Q|}$. Then, one of the following cases must hold.
\begin{enumerate}
\item There is at most one $\alpha\in\mathbb F$ for which there exists some $v\in L$ satisfying $v(Q)=x$\footnote{For a set $A=\{a_1,\ldots,a_{|A|}\}$, $v(A)$ denotes the sequence $(v(a_1),\ldots,v(a_{|A|}))$.} and $v(p)=\alpha$;
\item For every $\alpha\in\mathbb F$ there is an equal number of $v\in L$ for which $v(Q)=x$ and $v(p)=\alpha$.
\end{enumerate}
In particular, either no function (even randomized) of $v(Q)$ can predict $v(p)$ with probability more than $1/|\mathbb F|$, when $v\in L$ is randomly chosen uniformly, or $v(Q)$ always determines $v(p)$.
\end{nttheorem}
With that, we now prove Claim23\footnote{This proof is inspired by the proof of~\cite{ref23} of their Theorem1 and by a proof in~\cite{ref41} for a different claim.}.
\begin{proof}[Author proof for Claim23]
To show that $C$ is an SLR, we devise a recovering procedure $\operatorname{Rec}$ for it, based on the local corrector promised by it being an LCC. Let $D$\footnote{We make the slight assumption that $D^c(p)$ never directly queries $c(p)$. If however it does, then similarly $C$ can be shown to be a $(q,\tau,\varepsilon')$-SLR for $\tau=1/((\delta|P|/q)-1)$.} be such a local corrector. For every point $p\in P$, we construct a sequence of disjoint sets $Q_p^1,\ldots,Q_p^{m_p}\subseteq P$, where for every $i$, $c(Q_p^i)$ determines $c(p)$ while satisfying $|Q_p^i|\leq q$, and $m_p\geq\delta|P|/q$. On $p\in P$ and oracle access to $c\in C$, the procedure $\operatorname{Rec}^c(p)$ acts by uniformly choosing $i\in[m_p]$, querying $c(Q_p^i)$, and using it to deduce and output $c(p)$. The correctness of the result of $\operatorname{Rec}$ is immediate (since $\operatorname{Rec}$ always succeeds, $\varepsilon'=0$), and indeed the number of queries is no more than $q$. Since the sets are disjoint, the probability that a point is queried is no more than $\tau=q/(\delta|P|)$. It only remains to show how the assumed sets can be constructed, to conclude that $C$ is a $(q,\tau,\varepsilon')$-SLR, which we now turn to do.
For every $p\in P$, $Q_p^1,\ldots,Q_p^{m_p}$ are constructed as follows. Set $Q_p^0=\varnothing$. For $i=1,2,\ldots$, set $S_i=Q_p^0\cup\cdots\cup Q_p^{i-1}$. If $|S_i|>\delta|P|$, halt and set $m_p=i-1$. Otherwise, it holds that for every $c\in C$, for every modification of the coordinates in $S_i$ to some erroneous values, the decoder $D$ correctly outputs $c(p)$ with probability at least $1-\varepsilon$. An equivalent description of this case is the following: for every $c\in C$ and $z:S_i\to\mathbb F$, define $c_z\in\mathbb F^P$ such that for $r\notin S_i$, $c_z(r)=c(r)$, and for $r\in S_i$, $c_z(r)=z(r)$; the decoder $D$ chooses a set of queries $Q\subseteq P$, $|Q|\leq q$, according to a distribution and applies a function $f_Q$ on $c_z(Q)$; with probability at least $1-\varepsilon$, $f_Q(c_z(Q))=c(p)$. Since $Q$ is chosen in a manner independent of $c$ and $z$, one can verify that this implies that there exists some fixed $Q$ for which when $c\in C$ and $z:\mathbb F\to S_i$ are chosen randomly in a uniform manner, with probability at least $1-\varepsilon$ (this time over the choice of $c$ and $z$), $f_Q(c_z(Q))=c(p)$. Therefore, we can define another function $f'_Q$ that only gets $c(Q\setminus S_i)$, chooses $z$ at random, and outputs $f_Q(c_z(Q))$. If $c\in C$ is chosen uniformly at random, $f'_Q(Q\setminus S_i)=c(p)$ with probability at least $1-\varepsilon>1/|\mathbb F|$. By Fact24, this implies that $c(Q\setminus S_i)$ determines $c(p)$, for every $c\in C$. We therefore set $Q_p^i=Q\setminus S_i$, and proceed to next $i$. As this process only halts when $|S_i|>\delta|P|$, and for every $i$ $|S_i|\leq q(i-1)$, we have that indeed $m_p\geq\delta|P|/q$. Further note that by the choice of each $Q_p^i$, the sets $Q_p^1,\ldots,Q_p^{m_p}$ are disjoint, as required.
\end{proof}
\subsection*{4.1 Dual SLR and their induced SLR}
Our construction of SLR sets will be via constructing and analyzing sets which we call dual SLR sets. The SLR will then be induced from these dual SLR. We start by setting some notation. Let $P$ be a non-empty finite set and $\mathbb F$ a finite field. We make use of the standard notation $\mathbb F^P$ to denote the set of all functions $\{f:P\to\mathbb F\}$. Note that $\mathbb F^P$ has a natural $\mathbb F$-vector space structure where addition is point-wise, namely, for every $f,g\in\mathbb F^P$ and $a,b\in\mathbb F$ we have that $af+bg\in\mathbb F^P$ is defined by $(af+bg)(p)=af(p)+bg(p)$ for all $p\in P$. We consider the natural inner product map $\langle\cdot,\cdot\rangle:\mathbb F^P\times\mathbb F^P\to\mathbb F$ that is defined, for $f,g\in\mathbb F^P$, by $\langle f,g\rangle=\sum_{p\in P}f(p)g(p)$. Given $f\in\mathbb F^P$, we let $f^\perp=\{g\in\mathbb F^P\mid\langle f,g\rangle=0\}$. Note that $f^\perp$ is a linear subspace of $\mathbb F^P$. More generally, given a set $S\subseteq\mathbb F^P$ we define the linear subspace $S^\perp=\bigcap_{f\in S}f^\perp$. For $f\in\mathbb F^P$ we denote $|f|=|P\setminus f^{-1}(0)|$.
For the sake of readability, the field $\mathbb F$ and the set $P$ will be omitted from the notation that we are about the introduce in this section. Both will be clear from context. For $p\in P$ define $F_p=\{f\in\mathbb F^P\mid f(p)\ne0\}$. Informally, a dual SLR is a collection of distributions over $\mathbb F^P$, one for each point $p\in P$. The distribution $D_p$, that corresponds to $p$, outputs a function $g\in\mathbb F^P$. We think of $g$ as ``passing through'' $p$. We also allow $D_p$ to output a special ``failed symbol'' $\bot$ with some small probability. A dual SLR has the guarantee that $g$ does not pass through many other points, namely, $|g|$ is bounded, and that the dimension of all functions that can be sampled, when considering all distributions $D_p$, $p\in P$, is also bounded. Perhaps most importantly is the requirement that for every other fixed $r\in P$, the sampled $g\sim D_p$ is likely to have the property that $g\notin F_r$. Formally,
\begin{ntdefinition}[Author Definition25: Dual SLR]
Let $P$ be a set, $\mathbb F$ a field. Let $D=\{D_p\mid p\in P\}$ be a collection of distributions, where for each $p\in P$, $\operatorname{supp}(D_p)\subseteq F_p\cup\{\bot\}$. Denote $S=\bigcup_{p\in P}\operatorname{supp}(D_p)$. Let $L$ be a linear subspace of $\mathbb F^P$ such that $S\subseteq L\cup\{\bot\}$. The pair $(D,L)$ is said to be a $(q,\tau,\varepsilon,\rho)$-dual SLR on $\mathbb F^P$ provided the following holds:
\begin{enumerate}
\item $|g|\leq q$ for all $g\in S\setminus\{\bot\}$.
\item For every pair of distinct $p,r\in P$ (not necessarily distinct), it holds that $\Pr_{g\sim D_p}[g(r)\ne0\mid g\ne\bot]\leq\tau$.
\item For every $p\in P$, $\Pr[D_p=\bot]\leq\varepsilon$.
\item $\dim(L)\leq(1-\rho)|P|$.
\end{enumerate}
The linear subspace $L^\perp$ of $\mathbb F^P$ is called the induced SLR from the dual SLR $(D,L)$. The parameter $\tau$ of a dual-SLR is referred to as its smoothness.
\end{ntdefinition}
Let $(D,L)$ be a dual SLR. We turn to show that, as the name suggests, the induced SLR $L^\perp$ is indeed an SLR.
\begin{ntlemma}[Author Lemma26]
Let $P$ be a set, $\mathbb F$ a field, and let $(D,L)$ be $(q,\tau,\varepsilon,\rho)$-dual SLR on $\mathbb F^P$. Then the induced SLR $L^\perp$ is a $(q-1,\tau,\varepsilon)$-SLR. Furthermore, $L^\perp$ is linear and has rate $\rho$ or larger.
\end{ntlemma}
\begin{proof}
The moreover part readily follows since $L^\perp$ is a linear subspace of $\mathbb F^P$ and since $\dim(L^\perp)=|P|-\dim(L)\geq\rho|P|$. We describe a recovering procedure for $L^\perp$, namely, a randomized algorithm that is given an oracle access to $f\in L^\perp$ as well as a point $p\in P$ as input. The recovering procedure proceeds as follows:
\begin{enumerate}
\item Sample $g\sim D_p$. If $g=\bot$ return $\bot$; Otherwise,
\item Query $f$ on all points $Q=\{r\in P\setminus\{p\}\mid g(r)\ne0\}$.
\item Return $-\frac1{g(p)}\sum_{r\in Q}g(r)f(r)$.
\end{enumerate}
The query complexity of $\operatorname{Rec}$ is bounded by $q-1$ as $|Q|=|g|-1\leq q-1$. The probability that $\bot$ is returned is at most $\varepsilon$ by construction. We turn to prove that $\operatorname{Rec}^f(p)\in\{f(p),\bot\}$. By construction, $\operatorname{Rec}^f(p)=\bot$ if and only if $g=\bot$. Assume than that $g\ne\bot$, hence, $g\in\operatorname{supp}(D_p)\subseteq L$. As $f\in L^\perp$ we have that $0=\langle f,g\rangle$, and so
\[0=\sum_{r\in P}g(r)f(r)=g(p)f(p)+\sum_{r\in Q}g(r)f(r).\]
As $g\in\operatorname{supp}(D_p)\subseteq F_p$ we have $g(p)\ne0$, and so
\[f(p)=-\frac1{g(p)}\sum_{r\in Q}g(r)f(r)=\operatorname{Rec}^f(p).\]
To conclude the proof, we turn to analyze the smoothness of $\operatorname{Rec}$. First, note that, by construction, $f$ is never queried on $p$ itself. Consider then any $r\in P\setminus\{p\}$. Conditioned on $g\ne\bot$, the function $f$ is queried on $r$ if and only if $g(r)\ne0$. Thus,
\[\Pr[f(r)\text{ is queried}]=\Pr_{g\sim D_p}[g(r)\ne0\mid g\ne\bot]\leq\tau,\]
and the proof follows.
\end{proof}
We now show that the opposite holds as well, that any linear, non-adaptive, SLR induces a dual-SLR.
\begin{ntlemma}[Author Lemma27]
Let $P$ be a set, $\mathbb F$ a field, and let $C\subseteq\mathbb F^P$ be a linear non-adaptive $(q,\tau,\varepsilon)$-SLR with rate $\rho$. Then for some set $D$, $(D,C^\perp)$ is a $(q+1,\tau,\varepsilon,\rho)$-dual SLR.
\end{ntlemma}
\begin{proof}
Let $\operatorname{Rec}$ be a recovering procedure promised by $C$ being a SLR. Assume that $\operatorname{Rec}$ uses $R$ random bits. For every point $p\in P$, denote by $E_p$ the set of choices of the random bits $r\in\{0,1\}^R$ for which $\operatorname{Rec}^c(p)=\bot$ (for any $c\in C$). Note that $\mu(E_p\subseteq\{0,1\}^R)\leq\varepsilon$.
For any $p\in P$ and $r\in\{0,1\}^R$, $r\notin E_p$, denote by $Q_{p,r}\subseteq P$ the set of query locations which $\operatorname{Rec}(p)$ makes when $r$ is the choice of randomness. Define a function $f_{p,r}:C\to\mathbb F$ such that $f(c)$ is the output of $\operatorname{Rec}^c(p)$ when fixing its randomness to $r$ and note that $f_{p,r}(c)$ only depends on $\{c(w)\mid w\in Q_{p,r}\}$, and that $f_{p,r}(c)=c(p)$. Since $C$ is linear, one can easily verify that $f_{p,r}$ is a linear map. Therefore, for some $u_{p,r}\in\mathbb F^P$, $f_{p,r}(c)=\langle u_{p,r},c\rangle$ for every $c\in C$, where $u_{p,r}(w)=0$ if $w\notin Q_{p,r}$. We have that $c(p)=\langle u_{p,r},c\rangle$ for every $c\in C$. If we define a function $g_{p,r}\in\mathbb F^P$ such that
\[g_{p,r}(w)=\begin{cases}-1,&w=p;\\u_{p,r}(q),&w\ne p,\end{cases}\]
it follows that $g_{p,r}\in C^\perp$. Note that $|g_{\{p,r\}}|\leq q+1$.
For every $p\in P$, define $D_p$ to be the following distribution. To sample from $D_p$, draw $r\in\{0,1\}^R$ uniformly at random. If $r\in E_p$, output $\bot$; otherwise, output $g_{p,r}$. Set $D=\{D_p\mid p\in P\}$ and $L=C^\perp$. It follows trivially by the definitions that $(D,L)$ is a $(q+1,\tau,\varepsilon,\rho)$-dual SLR.
\end{proof}

\subsection*{4.2 Rate amplification for dual-induced SLR}
In this section we describe our first rate amplification procedure for SLR that are induced by dual SLR. Unlike the previous section, it will be more convenient to explicitly state within the notation the set $P$ over which we are working as we will be dealing with several such sets. The field $\mathbb F$, however, remains suppressed from the notation as it remains fixed in all SLR under consideration. We start by defining the following map of functions.
\setcounter{nttheorem}{27}
\begin{ntdefinition}[Author Definition28]
Let $P$ be a set and $\mathbb F$ a field. For an integer $\ell\geq1$ we define the map $\Phi:(\mathbb F^P)^\ell\to\mathbb F^{P^\ell}$ as follows. Let $g_1,\ldots,g_\ell\in\mathbb F^P$. The function $\Phi(g_1,\ldots,g_\ell):P^\ell\to\mathbb F$ is defined by
\[\Phi(g_1,\ldots,g_\ell)(p_1,\ldots,p_\ell)=\prod_{i=1}^\ell g_i(p_i)\]
for every $(p_1,\ldots,p_\ell)\in P^\ell$.
\end{ntdefinition}
Observe that $\Phi$ is multi-linear. Further, when $\ell=2$ and $g_1,g_2$ are viewed as vectors rather than functions, $\Phi$ is the outer product of the vectors.
\begin{ntdefinition}[Author Definition29]
Let $P$ be a set, $\mathbb F$ a field. Let $L^P$ be a linear subspace of $\mathbb F^P$. For an integer $\ell\geq1$, we define
\[L^{P^\ell}=\operatorname{Span}\{\Phi(g_1,\ldots,g_\ell)\mid g_1,\ldots,g_\ell\in L^P\}.\]
\end{ntdefinition}
\begin{ntclaim}[Author Claim30]
With the notation of Definition29, $\dim(L^{P^\ell})\leq\dim(L^P)^\ell$.
\end{ntclaim}
\begin{proof}
Let $B=\{g_1,\ldots,g_b\}$ be a basis for $L^P$, where $b=\dim(L^P)$. Define
\[B'=\{\Phi(h_1,\ldots,h_\ell)\mid(h_1,\ldots,h_\ell)\in B^\ell\}.\]
Observe that to prove the claim, it suffices to show that for every $f_1,\ldots,f_\ell\in L^P$ it holds that $\Phi(f_1,\ldots,f_\ell)\in\operatorname{Span}(B')$. As $f_1,\ldots,f_\ell\in L^P$, for every $i\in[\ell]$ we can write $f_i=\sum_{j=1}^b\lambda_{i,j}g_j$ with $\lambda_{i,j}\in\mathbb F$. We have that
\begin{align*}
\Phi(f_1,\ldots,f_\ell)&=\Phi\left(\sum_{j_1=1}^b\lambda_{1,j_1}g_{j_1},\ldots,\sum_{j_\ell=1}^b\lambda_{\ell,j_\ell}g_{j_\ell}\right)\\
&=\sum_{j_1,\ldots,j_\ell\in[b]}\left(\prod_{t=1}^\ell\lambda_{t,j_t}\right)\cdot\Phi(g_{j_1},\ldots,g_{j_\ell}),
\end{align*}
where the last equality follows by the multi-linearity of $\Phi$.
\end{proof}
\begin{ntdefinition}[Author Definition31]
Let $P$ be a set, $\mathbb F$ a field, and let $(D^P,L^P)$ be $(q,\tau,\varepsilon,\rho)$-dual SLR. Let $\ell\geq1$ be an integer. For $p\in P^\ell$ we define the distribution $D_p^{P^\ell}$ as follows. Write $p=(p_1,\ldots,p_\ell)$. To sample an element from $D_p^{P^\ell}$ proceed as follows:
\begin{enumerate}
\item Sample $g_1\sim D_{p_1}^P,\ldots,g_\ell\sim D_{p_\ell}^P$ independently.
\item If there exists $i\in[\ell]$ such that $g_i=\bot$, return $\bot$; Otherwise
\item Return $\Phi(g_1,\ldots,g_\ell)$.
\end{enumerate}
The collection of distributions $\{D_p^{P^\ell}\mid p\in P^\ell\}$ is denoted by $D^{P^\ell}$.
\end{ntdefinition}
We have the following lemma.
\begin{ntlemma}[Author Lemma32]
Let $P$ be a set, $\mathbb F$ a field, and let $(D^P,L^P)$ be a $(q,\tau,\varepsilon,\rho)$-dual SLR. Let $\ell\geq1$ be an integer and $D^{P^\ell}$ as in Definition31. Then, for every $p,r\in P^\ell$,
\[\Pr_{g\sim D_p^{P^\ell}}[g(r)\ne0\mid g\ne\bot]\leq\tau^{\operatorname{dist}(p,r)}.\]
\end{ntlemma}
\begin{proof}
Write $p=(p_1,\ldots,p_\ell)$, $r=(r_1,\ldots,r_\ell)$. By Definition31, conditioned on $g\ne\bot$ we have that $g=\Phi(g_1,\ldots,g_\ell)$ with $g_i\sim D_{p_i}^P$ for each $i\in[\ell]$ independently. Thus, $g(r)\ne0$ is the event
\[\Phi(g_1,\ldots,g_\ell)(r_1,\ldots,r_\ell)=\prod_{i=1}^\ell g_i(r_i)\ne0.\]
By the independence of $g_1,\ldots,g_\ell$, and since we are working over a field $\mathbb F$ (and so a product is nonzero if and only if each of the terms is nonzero), we get
\begin{equation}\Pr_{g\sim D_p^{P^\ell}}[g(r)\ne0\mid g\ne\bot]=\prod_{i=1}^\ell\Pr_{g_i\sim D_{p_i}^P}[g_i(r_i)\ne0\mid g_i\ne\bot].\tag{4.1}\end{equation}
Let $T=\{i\in[\ell]\mid p_i\ne r_i\}$. As $D^P$ is a $(q,\tau,\varepsilon,\rho)$-dual SLR, for each $i\in T$ it holds that
\[\Pr_{g_i\sim D_{p_i}^P}[g_i(r_i)\ne0\mid g_i\ne\bot]\leq\tau.\]
Substituting to Equation(4.1), we get
\[\Pr_{g\sim D_p^{P^\ell}}[g(r)\ne0\mid g\ne\bot]\leq\tau^{|T|},\]
which completes the proof.
\end{proof}
\begin{ntdefinition}[Author Definition33]
Let $P$ be a set, $\mathbb F$ a field, and let $(D^P,L^P)$ be a $(q,\tau,\varepsilon,\rho)$-dual SLR. For an integer $\ell\geq1$ let $L^{P^\ell},D^{P^\ell}$ be as in Definition29 and Definition31, respectively. We denote the pair $(D^{P^\ell},L^{P^\ell})$ by $(D^P,L^P)^\ell$.
\end{ntdefinition}
\begin{ntproposition}[Author Proposition34]
Let $P$ be a set, $\mathbb F$ a field, and let $(D^P,L^P)$ be a $(q,\tau,\varepsilon,\rho)$-dual SLR. Then, for every integer $\ell\geq1$ we have that $(D^P,L^P)^\ell$ is a $(q_\ell,\tau_\ell,\varepsilon_\ell,\rho_\ell)$-dual SLR, where
\[q_\ell\leq q^\ell,\quad\tau_\ell\leq\tau,\quad\varepsilon_\ell\leq\ell\varepsilon,\quad\rho_\ell\geq1-(1-\rho)^\ell.\]
\end{ntproposition}
\begin{proof}
First note that for every $p\in P^\ell$, the distribution $D_p^{P^\ell}$ is supported on $F_p^{P^\ell}\cup\{\bot\}$. Indeed, if we write $p=(p_1,\ldots,p_\ell)$ then, conditioned on $g\ne\bot$, we have that $g=\Phi(g_1,\ldots,g_\ell)$ where $g_i\in D_{p_i}^P$. Thus,
\[g(p)=\Phi(g_1,\ldots,g_\ell)(p_1,\ldots,p_\ell)=\prod_{i=1}^\ell g_i(p_i)\ne0.\]
Moreover, by Definition29,
\[\bigcup_{p\in P^\ell}\operatorname{supp}(D_p^{P^\ell})\subseteq L^{P^\ell}\cup\{\bot\}.\]
We turn to show that $q_\ell\leq q^\ell$. Let $p=(p_1,\ldots,p_\ell)\in P^\ell$ and consider any $g\in\operatorname{supp}(D_p^{P^\ell})$. By Definition31, if $g\ne\bot$ then $g=\Phi(g_1,\ldots,g_\ell)$ where $g_i\in\operatorname{supp}(D_{p_i}^P)\setminus\{\bot\}$. Now, for every $r=(r_1,\ldots,r_\ell)\in P^\ell$ we have that
\[g(r)\ne0\quad\Longleftrightarrow\quad\prod_{i=1}^\ell g_i(r_i)\ne0.\]
Since $\mathbb F$ is a field, the above is equivalent to $g_i(r_i)\ne0$ for all $i\in[\ell]$. Hence there are at most $q^\ell$ points $r\in P^\ell$ for which $g(r)\ne0$, and so $q_\ell\leq q^\ell$.
The bound on the smoothness readily follows by Lemma32. Indeed, consider any pair of distinct $p,r\in\mathbb F^{P^\ell}$. We have that $\operatorname{dist}(p,r)\geq1$ and so, by Lemma32,
\begin{equation}\Pr_{g\sim D_p^{P^\ell}}[g(r)\ne0\mid g\ne\bot]\leq\tau^{\operatorname{dist}(p,r)}\leq\tau.\tag{4.2}\end{equation}
To bound the probability that $\bot$ is returned, note that the event $D^{P^\ell}=\bot$ holds only if for some $i\in[\ell]$, $g_i=\bot$. Hence, by the union bound, $\Pr[D_p^{P^\ell}=\bot]\leq\ell\varepsilon$. We conclude the proof by bounding the dimension of $L^{P^\ell}$. By assumption, $\dim(L^P)\leq(1-\rho)|P|$. Claim30 then implies that
\[\dim(L^{P^\ell})\leq\dim(L^P)^\ell\leq((1-\rho)|P|)^\ell=(1-\rho)^\ell|P^\ell|.\]
\end{proof}
\paragraph{Discussion on the smoothness $\tau_\ell=\tau$}
The downside of the rate amplification procedure that was given in this section is that $\tau_\ell$ does not decrease with $\ell$ (which is bad as, recall, we wish $\tau$ to be small as, by Claim22, the distance $\delta$ of the resulted LCC is proportional to $1/\tau$). Indeed, with the notation of Proposition34, $\tau_\ell=\tau$. By examining the proof and Lemma32 one natural idea is to consider an SLR not over the entire set $P^\ell$ but on some subset of it which is a code with distance, say, $d>1$. This will indeed guarantee that for every two points $p,r$ we have $\operatorname{dist}(p,r)\geq d$ and so the bound in Equation(4.2) will be $\tau^d$ rather than $\tau$. While natural, this idea fails to yield better parameters as the rate-loss incurred by using a code (even an MDS) is larger than the improvement on the rate guaranteed via the rate amplification procedure.
In the next sections we give a more elaborate rate amplification procedure (that is based on the one that was given in this section) in which $\tau$ does decrease with $\ell$. Roughly, $\tau_\ell=(q\cdot\log|P|)^{\operatorname{poly}(\ell)}\tau^\ell$, and so there is a slight loss in the smoothness, which the reader should think as negligible. The query complexity $q_\ell$ as well as the rate $\rho_\ell$ and $\varepsilon_\ell$ are all slightly worse than those obtained in the above rate amplification procedure and so the two techniques are incomparable.

\subsection*{4.3 Distance-efficient rate amplification}
Let $P$ be a set, and $\mathcal R$ a partition of $P^2$. We denote the part containing $p$ by $[p]_{\mathcal R}$ or $[p]$ when $\mathcal R$ is clear from context. We call $(p)=[p]\setminus\{p\}$ the open class of $p$. For a set $A\subseteq P^2$ we let $(A)=\bigcup_{p\in A}(p)$. Given $p\in P$ we say that $\{p\}\times P\subseteq P^2$ is vertical line and $P\times\{p\}$ is a horizontal line. Horizontal and vertical lines are referred to as axis-parallel lines, and we denote the set of such lines by
\[\mathcal X=\bigcup_{p\in P}\{\{p\}\times P,P\times\{p\}\}.\]
For a point $p=(p_1,p_2)\in P^2$ we denote $S_p=(\{p_1\}\times P)\cup(P\times\{p_2\})\setminus\{p\}$. That is, $S_p$ is the set of points in $P^2$ of distance exactly $1$ from $p$. Key to our distance-efficient rate amplification procedure is a partition of the ``square'' $P^2$ with certain properties.
\setcounter{nttheorem}{34}
\begin{ntdefinition}[Author Definition35: Axis-evasive partitions]
Let $P$ be a set. A partition $\mathcal R$ of $P^2$ is said to be $(c,s)$-axis evasive if
\begin{enumerate}
\item For every $p\in P^2$, $|(p)|\leq c$.
\item For every $\ell,\ell'\in\mathcal X$ (possibly equal), $|\ell'\cap(\ell)|\leq s$.
\item For every $p\in P^2$ and $\ell\in\mathcal X$, $|[p]\cap\ell|\leq1$.
\end{enumerate}
\end{ntdefinition}
In Section5 we study such partitions. We prove their existence with certain parameters and give explicit constructions. In this section, however, we work with abstract axis-evasive partitions and analyze our rate amplification procedure with respect to the parameters $c,s$ of the axis-evasive partition as well as the number of parts which we typically denote by $t$.
\begin{ntclaim}[Author Claim36]
Let $p,p'\in P^2$ (possibly equal). Then, $|\{r\in S_p\mid(r)\cap S_{p'}\ne\varnothing\}|\leq4s$.
\end{ntclaim}
\begin{proof}
Note that each of $S_p,S_{p'}$ is a subset of the union of two axis-parallel lines. Thus, to prove the claim, it suffices to show that for every $\ell,\ell'\in\mathcal X$, not necessarily distinct,
\[|\{r\in\ell\mid(r)\cap\ell'\ne\varnothing\}|\leq s.\]
Let $r_1,\ldots,r_t\in\ell$ be such that $(r_i)\cap\ell'\ne\varnothing$. Note that for every distinct $i,j\in[t]$ it holds that $((r_i)\cap\ell')\cap((r_j)\cap\ell')=\varnothing$. Indeed, since $\mathcal R$ is a partition, if $((r_i)\cap\ell')\cap((r_j)\cap\ell')\ne\varnothing$ then $r_i\in[r_j]$, but this implies that $|\ell\cap[r_j]|\geq2$ in contradiction axis evasiveness. Thus,
\[R=\bigcup_{i=1}^t((r_i)\cap\ell')\]
is a disjoint union of size $t$. However, $R\subseteq(\ell)\cap\ell'$, and so $t\leq|R|\leq|(\ell)\cap\ell'|\leq s$.
\end{proof}
\begin{ntdefinition}[Author Definition37]
Let $P$ be a set, $\mathbb F$ a field. Let $\mathcal R$ be a $(c,s)$-axis evasive partition of $P^2$. For every $p\in P^2$ define the function $g_{[p]}:P^2\to\mathbb F$ as follows:
\[g_{[p]}(r)=\begin{cases}1,&r\in[p];\\0,&\text{otherwise}.\end{cases}\]
We define $L_{\mathcal R}=\{g_{[p]}\mid p\in P^2\}$.
\end{ntdefinition}
\begin{ntdefinition}[Author Definition38]
Let $P$ be a set, $\mathbb F$ a field. For $S\subseteq P$ define the function $\nu_S:P\to\mathbb F$ by
\[\nu_S(r)=\begin{cases}0,&r\in S;\\1,&\text{otherwise}.\end{cases}\]
For ease of readability, when $S$ is a singleton $S=\{p\}$, we write $\nu_p$ instead of $\nu_{\{p\}}$.
\end{ntdefinition}
With the notations and definitions above, we are ready to start developing our second rate amplification procedure. We start with the following.
\begin{ntdefinition}[Author Definition39]
Let $P$ be a set, $\mathbb F$ a field, and let $(D^P,L^P)$ be a $(q,\tau,\varepsilon,\rho)$-dual SLR. Let $L^{P^2}$ be as in Definition29. Let $\mathcal R$ be a $(c,s)$-axis evasive partition of $P^2$. We define for every $p\in P^2$ the distribution $(D_{\mathcal R}^{P^2})_p$ as follows. To sample $u$ from $(D_{\mathcal R}^{P^2})_p$:
\begin{enumerate}
\item Sample $g\sim D_p^{P^2}$.
\item If $g=\bot$ return $\bot$; Otherwise, denote $L=\{r\in S_p\mid g(r)\ne0\}$ and proceed as follows.
\item For every $r\in L$ and $w\in(r)$ sample $h_{r,w}\sim D_w^{P^2}$.
\item If there exist $r\in L$ and $w\in(r)$ such that either $h_{r,w}=\bot$ or $h_{r,w}(p)\ne0$ return $\bot$. Otherwise return
\begin{equation}u=g\nu_L+\sum_{r\in L}g(r)\sum_{w\in(r)}\frac{h_{r,w}\nu_w}{h_{r,w}(w)}.\tag{4.3}\end{equation}
\end{enumerate}
Note that, upon reaching Step(4), $u$ is well-defined as $h_{r,w}(w)\ne0$ for all $r\in L$ and $w\in(r)$. We denote the collection of distributions $\{(D_{\mathcal R}^{P^2})_p\mid p\in P^2\}$ by $D_{\mathcal R}^{P^2}$.
\end{ntdefinition}
We start by analyzing the function $u$ that is given by Equation(4.3) above.
\begin{ntclaim}[Author Claim40]
With the notation of Definition39, if $\bot$ is not returned then $u\in F_p$.
\end{ntclaim}
\begin{proof}
As $\bot$ was not returned, for every $r\in L$ and $w\in(r)$ it holds that $h_{r,w}\ne\bot$ and $h_{r,w}(p)=0$. Substituting to Equation(4.3), we get
\[u(p)=g(p)\nu_L(p)=g(p)\ne0,\]
where the second equality holds as $p\notin L$ and the last inequality follows since $g\in\operatorname{supp}(D_p^{P^2})\setminus\{\bot\}$.
\end{proof}
\begin{ntclaim}[Author Claim41]
With the notation of Definition39, if $\bot$ is not returned then $u\in L^{P^2}+L_{\mathcal R}$.
\end{ntclaim}
\begin{proof}
Take $f\in(L^{P^2}+L_{\mathcal R})^\perp$. To prove the claim, it suffices to show that $\langle u,f\rangle=0$. Indeed, this would imply $u\in((L^{P^2}+L_{\mathcal R})^\perp)^\perp=L^{P^2}+L_{\mathcal R}$. As $u\ne\bot$ we have that $g\ne\bot$. Note that
\[\langle g\nu_L,f\rangle=\langle g,f\rangle-\sum_{r\in L}g(r)f(r).\]
Since $g\in\operatorname{supp}(D_p^{P^2})$ we get that $g\in L^{P^2}$. However, $f\in(L^{P^2}+L_{\mathcal R})^\perp\subseteq(L^{P^2})^\perp$, implying $\langle g,f\rangle=0$. Thus,
\begin{equation}\langle g\nu_L,f\rangle=-\sum_{r\in L}g(r)f(r).\tag{4.4}\end{equation}
Now, fix $r\in L$ and $w\in(r)$. By Definition39, as $u\ne\bot$ we have that $h_{r,w}\ne\bot$ and so $h_{r,w}\in L^{P^2}$. However, by the above, $f\in(L^{P^2})^\perp$ and so $\langle h_{r,w},f\rangle=0$. Thus,
\[\langle h_{r,w}\nu_w,f\rangle=\langle h_{r,w},f\rangle-h_{r,w}(w)f(w)=-h_{r,w}(w)f(w).\]
Therefore, for every fixed $r\in L$ one has that
\begin{align}\left\langle\sum_{w\in(r)}\frac{h_{r,w}\nu_w}{h_{r,w}(w)},f\right\rangle&=\sum_{w\in(r)}\left\langle\frac{h_{r,w}\nu_w}{h_{r,w}(w)},f\right\rangle\notag\\&=\sum_{w\in(r)}\frac1{h_{r,w}(w)}\langle h_{r,w}\nu_w,f\rangle\notag\\&=-\sum_{w\in(r)}f(w).\tag{4.5}\end{align}
Now, $f\in(L^{P^2}+L_{\mathcal R})^\perp\subseteq(L_{\mathcal R})^\perp$ whereas $g_{[r]}\in L_{\mathcal R}$, and so
\[0=\langle f,g_{[r]}\rangle=\sum_{w\in[r]}f(w).\]
Substituting this to Equation(4.5), we get
\[\left\langle\sum_{w\in(r)}\frac{h_{r,w}\nu_w}{h_{r,w}(w)},f\right\rangle=f(r).\]
Therefore,
\begin{align*}\left\langle\sum_{r\in L}g(r)\sum_{w\in(r)}\frac{h_{r,w}\nu_w}{h_{r,w}(w)},f\right\rangle&=\sum_{r\in L}g(r)\left\langle\sum_{w\in(r)}\frac{h_{r,w}\nu_w}{h_{r,w}(w)},f\right\rangle\\&=\sum_{r\in L}g(r)f(r).\end{align*}
The above equation together with Equation(4.4) yield $\langle u,f\rangle=0$.
\end{proof}
\begin{ntclaim}[Author Claim42]
With the notation of Definition39, for every $p\in P^2$,
\[\Pr[(D_{\mathcal R}^{P^2})_p=\bot]\leq18csq\tau^2+2cq\varepsilon.\]
\end{ntclaim}
\begin{proof}
First, the probability that $g=\bot$ is bounded by $\varepsilon$. Similarly, the probability that for any specific $r\in L$ and $w\in(r)$, $h_{r,w}=\bot$ is bounded by $\varepsilon$. Thus, by the union bound, and since $|L|\leq2q-1$ and $|(r)|\leq c$, we have that expect with probability $(1+(2q-1)c)\varepsilon\leq2qc\varepsilon$, the sampling process above will result in $u\ne\bot$.
To complete the analysis, we turn to bound the probability that $h_{r,w}(p)=0$ for some $r\in L$ and $w\in(r)$. Let $L=\{r_1,\ldots,r_{|L|}\}$. While the random variables in $L$ may be dependent, marginally, it holds that for every $i\in[|L|]$ and every fixed $r\in S_p$, $\Pr[r_i=r]\leq\tau$. With this notation, by Definition39, $(D_{\mathcal R}^{P^2})_p=\bot$ only if there exist $i\in[|L|]$ and $w\in(r_i)$ such that $h_{r_i,w}(p)\ne0$.
For a fixed $r\in S_p$ define the event $E_r$ in which there exists $w\in(r)$ such that $h_{r,w}(p)\ne0$, (when conditioned on $h_{r,w}\ne\bot$). Note that this event is with respect to the randomness of sampling $h_r=\{h_{r,w}\mid w\in(r)\}$ whereas $r$ is fixed. By the union bound,
\[\Pr_{h_r}[E_r]\leq\sum_{w\in(r)}\Pr_{h_{r,w}}[h_{r,w}(p)\ne0\mid h_{r,w}\ne\bot].\]
Observe first that $w\ne p$. Indeed, as $r\in S_p$, both $r$ and $p$ are on some common axis-parallel line $\ell\in\mathcal X$. Thus, $w=p$ would imply $|[r]\cap\ell|\geq2$ which stands in contradiction to the definition of axis-evasiveness. Consider $w\in(r)\setminus S_p$. As $w\ne p$ we have that $\operatorname{dist}(w,p)=2$. By Lemma32, as $h_{r,w}\sim D_w^{P^2}$ we have that
\[\Pr_{h_{r,w}}[h_{r,w}(p)\ne0\mid h_{r,w}\ne\bot]\leq\tau^2.\]
If, on the other hand, $w\in(r)\cap S_p$ then $\operatorname{dist}(w,p)=1$, and Lemma32 then implies that
\[\Pr_{h_{r,w}}[h_{r,w}(p)\ne0\mid h_{r,w}\ne\bot]\leq\tau.\]
As $|(r)|\leq c$ we conclude that
\[\Pr_{h_r}[E_r]\leq c\tau^2+\tau|(r)\cap S_p|.\]
Fix $i\in[|L|]$ and consider the random variable $r_i$. The above equation, together with $|(r_i)|\leq c$, yields
\begin{align}\Pr_{r_i,h_{r_i}}[E_{r_i}]&\leq\Pr_{r_i,h_{r_i}}[E_{r_i}\mid(r_i)\cap S_p=\varnothing]\notag\\&\quad+\Pr_{r_i,h_{r_i}}[E_{r_i}\mid(r_i)\cap S_p\ne\varnothing]\Pr_{r_i}[(r_i)\cap S_p\ne\varnothing]\notag\\&\leq c\tau^2+(c\tau^2+c\tau)\Pr_{r_i}[(r_i)\cap S_p\ne\varnothing].\tag{4.6}\end{align}
Consider now the set $B=\{r\in S_p\mid(r)\cap S_p\ne\varnothing\}$. As $\mathcal R$ is $(c,s)$-axis evasive, Claim36 implies that $|B|\leq4s$, and so
\[\Pr_{r_i}[(r_i)\cap S_p\ne\varnothing]=\Pr[r_i\in B]\leq4s\tau.\]
Substituting to Equation(4.6), we get $\Pr[E_i]\leq9cs\tau^2$. The proof then follows by taking the union bound over all $i\in[|L|]$ as, indeed, $|L|=2q-1$.
\end{proof}
\begin{ntclaim}[Author Claim43]
With the notation of Definition39, for every pair of distinct $p,r\in P^2$,
\[\Pr_{u\sim(D_{\mathcal R}^{P^2})_p}[u(r)\ne0\mid u\ne\bot]\leq10csq\tau^2.\]
\end{ntclaim}
\begin{proof}
By Equation(4.3), $u$ is a linear combination of the (sampled) functions $g\nu_L,\{h_{r,w}\nu_w\}$. To prove the claim, we will show that, with high probability, all these functions evaluate to $0$ at the point $r$, implying $u(r)=0$. We start by bounding $\Pr[(g\nu_L)(r)\ne0]$. To this end, consider two cases. First, if $r\in P^2\setminus S_p$ then, as $L\subseteq S_p$, we have that $\nu_L(r)=1$ and so in such case
\begin{equation}\Pr[(g\nu_L)(r)\ne0]=\Pr[g(r)\ne0]\leq\tau^2,\tag{4.7}\end{equation}
where the last inequality follows by Lemma32 and since $\operatorname{dist}(r,p)=2$ per our assumption $r\notin S_p$ and since $r\ne p$. If, on the other hand, $r\in S_p$ then, by the definition of $L$,
\[g(r)\ne0\Longrightarrow r\in L\Longrightarrow\nu_L(r)=0,\]
and so in this case $(g\nu_L)(r)=0$.
Let $L=\{r_1,\ldots,r_{|L|}\}$. Consider a fixed $i\in[|L|]$ and denote $(r_i)=\{w_{i,1},\ldots,w_{i,b}\}$, where $b\leq c$. Fix $j\in[b]$. We turn to bound $\Pr[(h_{r_i,w_{i,j}}\nu_{w_{i,j}})(r)\ne0]$. First note that
\begin{equation}\Pr[(h_{r_i,w_{i,j}}\nu_{w_{i,j}})(r)\ne0\mid(r_i)\cap S_r=\varnothing]\leq\tau^2.\tag{4.8}\end{equation}
Indeed, conditioned on the event $(r_i)\cap S_r=\varnothing$, either $w_{i,j}=r$ or $\operatorname{dist}(w_{i,j},r)=2$. In the first case,
\[(h_{r_i,w_{i,j}}\nu_{w_{i,j}})(r)=h_{r_i,r}(r)\nu_r(r)=0.\]
In the second case, the bound follows by Lemma32. Second, note that
\begin{equation}\Pr[(h_{r_i,w_{i,j}}\nu_{w_{i,j}})(r)\ne0\mid(r_i)\cap S_r\ne\varnothing]\leq\tau.\tag{4.9}\end{equation}
Indeed, as before, we may only consider the case $r\ne w_{i,j}$ and then observe that $\operatorname{dist}(r,w_{i,j})=1$ and invoke Lemma32. Now, let $B=\{v\in S_p\mid(v)\cap S_r\ne\varnothing\}$. By Claim36, and since $\mathcal R$ is $(c,s)$-axis evasive, $|B|\leq4s$. Recall that $\Pr[r_i=v]\leq\tau$ for every fixed $v\in S_p$, and so
\begin{equation}\Pr[(r_i)\cap S_r\ne\varnothing]=\Pr[r_i\in B]\leq4s\tau.\tag{4.10}\end{equation}
By Equations(4.8),(4.9),(4.10) we get
\[\Pr[(h_{r_i,w_{i,j}}\nu_{w_{i,j}})(r)\ne0]\leq\tau^2+4s\tau^2\leq5s\tau^2.\]
The proof then follows by the union bound over all $i\in[|L|]$ and $j\in[|(w_i)|]$.
\end{proof}
\begin{ntdefinition}[Author Definition44]
Let $P$ be a set, $\mathbb F$ a field, and let $(D^P,L^P)$ be a $(q,\tau,\varepsilon,\rho)$-dual SLR. Let $L^{P^2}$ be as in Definition29. Let $\mathcal R$ be a $(c,s)$-axis evasive partition of $P^2$ and let $D_{\mathcal R}^{P^2}$ be as in Definition39. We denote by $(D^P,L^P)^2_{\mathcal R}$ the pair $(D_{\mathcal R}^{P^2},L^{P^2}+L_{\mathcal R})$.
\end{ntdefinition}
\begin{ntproposition}[Author Proposition45]
Let $P$ be a set, $\mathbb F$ a field, and let $(D^P,L^P)$ be a $(q,\tau,\varepsilon,\rho)$-dual SLR. Let $\mathcal R$ be a $(c,s)$-axis evasive partition of $P^2$ that consists of $t$ parts. Then, $(D^P,L^P)^2_{\mathcal R}$ is a $(q_{\mathcal R},\tau_{\mathcal R},\varepsilon_{\mathcal R},\rho_{\mathcal R})$-dual SLR with
\begin{align*}q_{\mathcal R}&\leq2cq^3\\\tau_{\mathcal R}&\leq10csq\tau^2\\\varepsilon_{\mathcal R}&\leq18csq\tau^2+2cq\varepsilon\\\rho_{\mathcal R}&\geq1-(1-\rho)^2-\frac t{|P|^2}.\end{align*}
\end{ntproposition}
\begin{proof}
Claim40 implies that for every $p\in P^2$, $\operatorname{supp}((D_{\mathcal R}^{P^2})_p)\subseteq F_p\cup\{\bot\}$. To bound $q_{\mathcal R}$, note that by Equation(4.3),
\[|u|\leq|g\nu_L|+\sum_{r\in L}\sum_{w\in(r)}|h_{r,w}\nu_w|.\]
Now, $|g\nu_L|\leq|g|\leq q^2$ and $|h_{r,w}\nu_w|\leq|h_{r,w}|\leq q^2$. Hence, $|u|\leq q^2+|L|cq^2\leq2cq^3$. The stated bounds on $\tau_{\mathcal R}$ and $\varepsilon_{\mathcal R}$ readily follows by Claim43 and Claim42, respectively. As for the rate, we have that
\begin{align*}\dim(L^{P^2}+L_{\mathcal R})&\leq\dim(L^{P^2})+\dim(L_{\mathcal R})\\&\leq(1-\rho)^2|P|^2+t,\end{align*}
where the second inequality follows by Proposition34 and since $\mathcal R$ consists of $t$ parts, implying $|L_{\mathcal R}|=t$.
\end{proof}

\subsection*{4.4 Proofs of Theorem3 and Corollary4}
With the machinery developed in the previous section, and using in a black-box manner, the construction of axis-evasive partitions we obtain in Section5, we are finally ready to prove Theorem3 and Corollary4. We start by giving a more formal statement of Corollary4 (Theorem46, stated at the start of this item).
\paragraph{A remark regarding the error}
Note that there is another implicit constraint on $\rho_{\mathrm{in}}$ and $\alpha$ that originates from the error. Indeed, to make the result non-trivial, one must have $\varepsilon_{\mathrm{out}}<1$ which, in turn, implies some bound on $\ell$ and then, through Equation(4.12), a constraint on $\rho_{\mathrm{in}}$ and $\alpha$. However, if that turns out to be a problem for the regime of parameters one is interested in, the probability to output $\bot$ can be reduced by repetition. Thus, by performing an alternating sequence of such error (or failure) reductions and rate amplifications, one can resolve this issue. Note that unlike for LDC, the error reduction has no cost in query complexity, and it certainly has no effect on the smoothness nor on the rate. It does, however, effects the running-time.
As mentioned above, our proof relies on an explicit axis-evasive partition that we construct in Section5. Formally,
\setcounter{nttheorem}{46}
\begin{nttheorem}[Author Theorem47]
Let $P$ be a set of size $q$, where $q$ is an odd prime power. Let $c$ be an even integer such that $c+1\mid q+1$, and $c\leq\sqrt q/10$. Then, there exists a $(c,4c^2)$-axis evasive partition of $P^2$ with at most $2q^2/(c+1)$ parts.
\end{nttheorem}
Our proof of Theorem46 is done by applying Proposition45 several times, iteratively, where in each iteration we square the size of the set $P$ obtained by the previous iterative step. Note, however, that Theorem47 requires the set size $|P|$ to be an odd prime power $q$ with the property that $c+1\mid q+1$. It is best to choose $c$ the same in all applications of Proposition45. However, note that if we start an iteration with a set of size $q$ and so end the iteration with a set of size $q^2$ then the condition will fail to hold at the beginning of the following iteration. Indeed if $c+1\mid q+1$ then $q\equiv-1\pmod{c+1}$ and so $q^2\equiv1\pmod{c+1}$. To overcome this technicality, we do not work with the set obtained by the previous iteration as is. Instead, we find a prime--not much smaller than $q^2$--that has the desired residue $-1$ modulo $c+1$. To this end we rely on the Siegel--Walfisz Theorem~\cite{ref34,ref37} which refines Dirichlet's theorem on primes in arithmetic progressions. The state the Siegel--Walfisz Theorem we set some notation. For an integer $m\geq1$, we denote Euler's totient function, that counts the positive integers up to $m$ that are relatively prime to $m$, by $\varphi(m)$. For integers $n,m,r$, we denote the number of (positive) primes less than or equal to $n$ which are congruent to $r$ modulo $m$ by $\pi(n;m,r)$. The Eulerian logarithmic integral is given by $\operatorname{Li}(x)=\int_2^x dt/\ln t$.
\begin{nttheorem}[Author Theorem48, {\cite{ref34,ref37}}]
For every constant $e\geq1$ there exists a constant $c=c(e)$ such that the following holds. Let $n,m,r$ be positive integers such that $m\leq(\log n)^e$, and $m,r$ coprimes. Then,
\[\left|\pi(n;m,r)-\frac{\operatorname{Li}(n)}{\varphi(m)}\right|=O\left(n\cdot2^{-c\sqrt{\log n}}\right).\]
\end{nttheorem}
We have the following straightforward corollary.
\begin{nttheorem}[Author Corollary49]
For every constant $e\geq1$ there exist constants $c=c(e),n_0=n_0(e)$ such that the following holds. Let $m,r$ be coprime integers, $m>0$. Let $n\geq n_0$ be an integer such that $m\leq(\log n)^e$. Then, there exists a prime $p\in[n-\Delta,n]$, where $\Delta=cn/\log n$, such that $p\equiv r\pmod m$.
\end{nttheorem}
\begin{proof}
To prove the corollary, it suffices to show that $\pi(n;m,r)>\pi(n-\Delta;m,r)$. By Theorem48, there exist constants $n_0,c'$ such that for every $n\geq n_0$,
\[\left|\pi(n;m,r)-\frac{\operatorname{Li}(n)}{\varphi(m)}\right|\leq c'n\cdot2^{-c\sqrt{\log n}}.\]
Thus, it suffices to show that
\[\frac{\operatorname{Li}(n)}{\varphi(m)}-c'n\cdot2^{-c\sqrt{\log n}}>\frac{\operatorname{Li}(n-\Delta)}{\varphi(m)}+c'(n-\Delta)\cdot2^{-c\sqrt{\log(n-\Delta)}}.\]
As we may assume that $\Delta\leq n/2$, it suffices to prove that
\begin{equation}\operatorname{Li}(n)-\operatorname{Li}(n-\Delta)\geq2c'\varphi(m)n\cdot2^{-c\sqrt{\log(n/2)}}.\tag{4.13}\end{equation}
It is well-known that
\[\operatorname{Li}(x)=c_1+\frac x{\ln x}+O\left(\frac x{\ln^2x}\right),\]
where $c_1=\int_{t=0}^2dt/\ln t$ is some constant. Therefore,
\[\operatorname{Li}(n)-\operatorname{Li}(n-\Delta)\geq\frac\Delta{\ln(n/2)}-\frac{c''n}{\ln^2n}.\]
for some constant $c''$. By our assumption on $\Delta$ we can choose the parameter $c$ in the definition of $\Delta$ such that the right hand side is bounded below by $n/\ln^2n$. The proof then follows by Equation(4.13) and noting that $\varphi(m)\leq m\leq(\log n)^e=o(2^{-c\sqrt{\log(n/2)}})$.
\end{proof}
We turn to formally define and analyze the operation of projecting a dual SLR over $\mathbb F^P$ on a (large) subset of $P$.
\begin{ntdefinition}[Author Definition50]
Let $P$ a set and $P'\subseteq P$. Let $p'\in P'$ and $D$ be a distribution with $\operatorname{supp}(D)\subseteq F_{p'}\cup\{\bot\}$. We define the $D|_{P'}$ as follows: To sample from $D|_{P'}$, sample $f\sim D$. If $f=\bot$, output $\bot$; if $f\in F_{p'}$, output $f|_{P'}$. We refer to $D|_{P'}$ as the distribution $D$ projected to $P'$.
\end{ntdefinition}
\begin{ntdefinition}[Author Definition51]
Let $P$ be a set, $\mathbb F$ a field. Let $D=\{D_p\mid p\in P\}$ be a collection of distributions, where for each $p\in P$, $\operatorname{supp}(D_p)\subseteq F_p\cup\{\bot\}$. Let $P'\subseteq P$. We define $D|_{P'}$ to be the collection $D$ projected to $P'$, that is, $D|_{P'}=\{D_{p'}|_{P'}\mid p'\in P'\}$.
\end{ntdefinition}
\begin{ntdefinition}[Author Definition52]
Let $P$ be a set, $\mathbb F$ a field and let $L$ be a linear subspace of $\mathbb F^P$. Let $P'\subseteq P$. We denote by $L|_{P'}$ the linear subspace $L$ projected to $P'$, namely, $L|_{P'}=\{f|_{P'}\mid f\in L\}$.
\end{ntdefinition}
\begin{ntclaim}[Author Claim53]
Let $P$ be a set, $\mathbb F$ a field, $(D,L)$ a $(q,\tau,\varepsilon,\rho)$-dual SLR over $\mathbb F^P$, and let $P'\subseteq P$. Then, $(D|_{P'},L_{P'})$ is a $(q,\tau,\varepsilon,\rho')$-dual SLR over $\mathbb F^{P'}$, where $\rho'=1-\frac{|P|}{|P'|}(1-\rho)$.
\end{ntclaim}
\begin{proof}
That the smoothness $\tau$, as well as $q$ and $\varepsilon$, all stay the same after projecting the dual SLR to $P'$, follows immediately from the definitions. The assertion regarding the rate of the induced SLR, $\rho'$, readily follows as we have that
\[\dim(L|_{P'})\leq\dim(L)\leq(1-\rho)|P|=\left(1-\left(1-\frac{|P|}{|P'|}(1-\rho)\right)\right)|P'|.\]
\end{proof}
\begin{ntclaim}[Author Claim54]
There exists a universal constant $m_0$ such that the following holds. Let $P$ be a set of size $m\geq m_0$. Let $\mathbb F$ be a field, and let $(D^P,L^P)$ be a $(q_{\mathrm{in}},\tau_{\mathrm{in}},\varepsilon_{\mathrm{in}},\rho_{\mathrm{in}})$-dual SLR over $\mathbb F^P$. Let $c\leq\log m$ be an integer. Then, there exists a set $P'$ of size
\[|P'|\geq\left(1-O\left(\frac1{\log m}\right)\right)m^2,\]
and a $(q_{\mathrm{out}},\tau_{\mathrm{out}},\varepsilon_{\mathrm{out}},\rho_{\mathrm{out}})$-dual SLR $(D^{P'},L^{P'})$ over $\mathbb F^{P'}$, where
\begin{align*}q_{\mathrm{out}}&\leq2cq_{\mathrm{in}}^3,\\\tau_{\mathrm{out}}&\leq40c^3q_{\mathrm{in}}\tau_{\mathrm{in}}^2,\\\varepsilon_{\mathrm{out}}&\leq80c^3q_{\mathrm{in}}(\tau_{\mathrm{in}}^2+\varepsilon_{\mathrm{in}}),\\\rho_{\mathrm{out}}&\geq1-(1-\rho_{\mathrm{in}})^2-O(1/c).\end{align*}
\end{ntclaim}
\begin{proof}
By Corollary49 applied with $n,m,r$ in the notation of Corollary49 set to $m,c+1,-1$ in the notation of this claim, respectively, there exists some prime $p\leq m$ such that $m-p=O(m/\log m)$, and $c+1\mid p+1$. Take $P'$ to be an arbitrary subset of $P$ of size $p$. By Claim53, $(D|_{P'},L|_{P'})$ is a $(q_{\mathrm{in}},\tau_{\mathrm{in}},\varepsilon_{\mathrm{in}},\rho')$-dual SLR on $P'$, where
\[\rho'=1-\frac mp(1-\rho_{\mathrm{in}})\geq\rho_{\mathrm{in}}-O\left(\frac1{\log m}\right).\]
By Theorem47 applied to $P'$, which observe is indeed applicable as $c+1\mid p+1$, there exists an explicit $(c,4c^2)$-axis evasive partition $\mathcal R$ of $(P')^2$ with at most $t=2p^2/(c+1)$ parts. With that partition, we can now apply Proposition45 to $(D|_{P'},L|_{P'})$ and get that $(D|_{P'},L|_{P'})^2_{\mathcal R}$ is a $(q_{\mathrm{out}},\tau_{\mathrm{out}},\varepsilon_{\mathrm{out}},\rho_{\mathrm{out}})$-dual SLR with the stated parameters. Note that the assertion regarding the rate follows as $c\leq\log m$,
\end{proof}
The following proposition is a more formal and accurate restatement of Theorem3.
\begin{ntproposition}[Author Proposition55]
There exist universal constants $0<c'<1$ and $c'',m',\ell'\geq1$ such that the following holds. Let $P$ be a set of size $m\geq m'$. Let $\mathbb F$ be a field, and let $(D^P,L^P)$ be a $(q_{\mathrm{in}},\tau_{\mathrm{in}},\varepsilon_{\mathrm{in}},\rho_{\mathrm{in}})$-dual SLR over $\mathbb F^P$. Let $\ell=2^r$ for an integer $r\geq1$, and assume that $\ell\geq\ell'$. Let $c$ be an integer such that $c''\ell^2\leq c\leq c'\log m$. Then, there exists a set $P_\ell$ of size $m^\ell/2\leq|P_\ell|\leq m^\ell$, and a $(q_\ell,\tau_\ell,\varepsilon_\ell,\rho_\ell)$-dual SLR $(D^{P_\ell},L^{P_\ell})$ over $\mathbb F^{P_\ell}$, where
\begin{align*}q_\ell&\leq(2cq_{\mathrm{in}})^{\ell^{\log3}},\\\tau_\ell&=O((c^3q_{\mathrm{in}})^{\ell^{\log3}})\cdot\tau_{\mathrm{in}}^\ell,\\\varepsilon_\ell&\leq O((c^4q_{\mathrm{in}})^{\ell^{\log3}})\cdot(\tau_{\mathrm{in}}+\varepsilon_{\mathrm{in}}),\\\rho_\ell&\geq1-(1-\rho_{\mathrm{in}})^\ell-O\left(\frac{\ell^2}c\right),\end{align*}
where, recall, the $\log$ function is taken base $2$.
\end{ntproposition}
\begin{proof}
We construct a sequence of $(q_t,\tau_t,\varepsilon_t,\rho_t)$-dual SLR $(D^{P_t},L^{P_t})$ for $t=0,1,\ldots,r=\log\ell$, and show that the last dual-SLR in the sequence has the stated parameters. The first dual-SLR, $(D^{P_0},L^{P_0})$, is taken to be the $(q_{\mathrm{in}},\tau_{\mathrm{in}},\varepsilon_{\mathrm{in}},\rho_{\mathrm{in}})$-dual SLR $(D^P,L^P)$ that is given by the hypothesis of the proposition. After constructing $(D^{P_t},L^{P_t})$, we obtain $(D^{P_{t+1}},L^{P_{t+1}})$ by applying Claim54 to $(D^{P_t},L^{P_t})$ with the parameter $c$ taken to be $c$ from the statement of this proposition. Note that, as required by the claim, $c\leq\log m$. Note that, by taking $m'$ to be a large enough constant, all other dual SLR in the sequence will have $|P_t|\geq m$ as well, and so we can apply Claim54 to them. Denote $m_t=|P_t|$. We begin by bounding $m_t$ from below. Indeed, by Claim54, and using that $1-x\geq e^{-2x}$ for $x\leq1/2$, we can pick the constant $c''$ such that
\[m_t\geq e^{-c''/\log m_{t-1}}m_{t-1}^2\geq e^{-c''/\log m_0}m_{t-1}^2,\]
where the last inequality follows as, for a large enough constant $m'$, the sequence $(m_t)_t$ is monotone increasing. We invoke Claim82 with $a=e^{c''/(2\log m_0)}$ and $b=2$ to conclude that
\[m_t\geq m_0^{2^t}e^{-c''2^t/\log m_0}\geq\tfrac12m_0^{2^t},\]
where the last inequality follows as $t\leq r=\log\ell$ and, recall, we take $\ell\leq c'\log m$ for a sufficiently small constant $c'>0$. In particular, $m_r\geq m^\ell/2$ as stated.
By Claim54, for every $t\geq1$ we have $q_t\leq2cq_{t-1}^3$. It is straightforward to prove by that
\begin{equation}q_t\leq(2cq_{\mathrm{in}})^{3^t},\tag{4.14}\end{equation}
which readily implies the assertion regarding the query complexity. We turn to analyze the rate. Denote $\beta_t=1-\rho_t$. Claim54 implies that $\beta_t\leq\beta_{t-1}^2+c'''/c$, for some constant $c'''>0$. By induction on $t$, we get that $\beta_t\leq\beta_0^{2^t}+c'''4^t/c$. Indeed, the base case $t=0$ is obvious. Now, by the induction hypothesis,
\[\beta_t\leq\beta_{t-1}^2+\frac{c'''}c\leq\left(\beta_0^{2^{t-1}}+\frac{c'''}c4^{t-1}\right)^2+\frac{c'''}c.\]
One can easily verify that the right hand side is bounded above by the desired bound $\beta_0^{2^t}+c'''4^t/c$ provided that $2^tc'''/c\leq1$. As $t\leq r$ and $2^r=\ell$, the latter inequality follows assuming $c'''\ell\leq c$. As we assume $c\geq c''\ell^2$, it suffices to choose $\ell'$ from the statement of the proposition to be a constant larger than the constant $c'''/c''$. We conclude that,
\[\beta_r\leq\beta_0^\ell+O\left(\frac{4^r}c\right)=\beta_0^\ell+O\left(\frac{\ell^2}c\right),\]
which implies the assertion regarding the rate.
As for the smoothness, by Claim54, and using Equation(4.14), we have that
\[\tau_t\leq40c^3q_{t-1}\tau_{t-1}^2\leq40c^3(2cq_{\mathrm{in}})^{3^{t-1}}\tau_{t-1}^2,\]
from which it is easy to verify that
\[\tau_t\leq(40c^3)^{2^t}(2cq_{\mathrm{in}})^{3^t}\tau_0^{2^t},\]
and the assertion regarding the smoothness readily follows. Last is the error which we leave to the reader to verify.
\end{proof}
We can now easily deduce Theorem46.
\begin{proof}[Author proof of Theorem46]
The proof readily follows from Proposition55 by taking $\ell$ as defined in Equation(4.12), and with $c$ in the notation of Proposition55 taken to be $c=\Theta(\ell^2/\alpha)$. Note that this choice of parameters satisfies the hypothesis of Proposition55 as indeed implied by Equation(4.11) and by taking $c'$ to be a sufficiently large constant. It is easy to verify that the rate is $1-\alpha$ with our choice of $c,\ell$, and the remaining assertions readily follow by Proposition55.
\end{proof}
\setcounter{nttheorem}{81}
\begin{ntclaim}[Author Claim82, used above; p1:48]
Let $(n_t)_{t\in\mathbb N}$ be a sequence of positive integers such that $n_t\geq(n_{t-1}/a)^b$ for some $a,b>1$. Then, for every $t\geq1$ we have that $n_t\geq(n_0/a^{h(b,t)})^{b^t}$, where $h(b,t)=\sum_{i=0}^{t-1}1/b^i$.
\end{ntclaim}
One can prove the following easy claim by induction (author's supplied justification for Claim82).

\section*{5 Axis-evasive partitions: 5.2 Explicit constructions}
In this section we give explicit constructions of axis-evasive partitions (see Definition35). Our constructions are based on quadratic field extensions. We identify a set $P$ of size $q$--a prime power--with the finite field $\mathbb F_q$ in an arbitrary manner, namely, by using an arbitrary bijection which, for ease of readability, we do not make explicit in the notation. We start by giving some basic background on finite fields.
Let $h(x)\in\mathbb F_q[x]$ be a degree $2$ irreducible monic polynomial. It is a well-known fact that $\mathbb F_q[x]/\langle h(x)\rangle$ is a field of size $q^2$ which we denote, somewhat less informatively, by $\mathbb F_{q^2}$. Note that there exists $\alpha\in\mathbb F_{q^2}$ such that $h(\alpha)=0$ (indeed, take $\alpha=x+\langle h(x)\rangle$). Since $h$ is irreducible over $\mathbb F_q$ and has degree $2$, we can write every element of $\mathbb F_{q^2}$ in the form $a+\alpha b$, where $a,b\in\mathbb F_q$, in a unique manner. That is, we can identify in the set-theoretic level, $\mathbb F_{q^2}$ with $\mathbb F_q+\alpha\mathbb F_q$. Using this identification, we identify $P^2$ with $\mathbb F_{q^2}$ in the natural way, namely, a point $(a,b)\in P^2$ is identified with $a+\alpha b$ in $\mathbb F_{q^2}$. Note that, with this identification, the horizontal lines in $P^2$ are of the form $b\alpha+\mathbb F_q$ where $b\in\mathbb F_q$ can be thought of as the fixed height of the line. Similarly, the vertical lines are given by $b+\alpha\mathbb F_q$. Given $\delta\in\mathbb F_{q^2}\setminus\{0\}$, we say that $\ell_\delta=\delta\mathbb F_q\subseteq\mathbb F_{q^2}$ is the line through the origin with slope $\delta$.
Our construction of exis-evasive partitions is based on an equivalence relation that we are about to define. The partition is then obtained by considering the respective equivalence classes. We begin the construction by ignoring the ``origin'' $0\in\mathbb F_{q^2}$ and work only with $\mathbb F_{q^2}\setminus\{0\}$. Note that this is the set of invertible elements of $\mathbb F_{q^2}$ which has a group structure under the field multiplication. When referring to this multiplicative group we write $(\mathbb F_{q^2})^\times$.
Let $\beta\in(\mathbb F_{q^2})^\times$. Denote by $o(\beta)$ the order of $\beta$ in the multiplicative group $(\mathbb F_{q^2})^\times$. It will be convenient to denote $c=o(\beta)-1$. We define an equivalence relation on $(\mathbb F_{q^2})^\times$, parameterized by $\beta$, as follows: For $\gamma,\delta\in(\mathbb F_{q^2})^\times$
\begin{equation}\gamma\sim\delta\quad\Longleftrightarrow\quad\gamma\delta^{-1}\in\langle\beta\rangle,\tag{5.1}\end{equation}
where $\langle\beta\rangle$ is the subgroup of $(\mathbb F_{q^2})^\times$ that is generated by $\beta$. Observe that this is an equivalence relation. Indeed, the classes are the different cosets, that is, the elements of the quotient group $(\mathbb F_{q^2})^\times/\langle\beta\rangle$. For completeness, we quickly prove that this is an equivalence relation: as $1\in\langle\beta\rangle$, we have that $\gamma\sim\gamma$. Secondly, if $\gamma\delta^{-1}\in\langle\beta\rangle$ then $\delta\gamma^{-1}\in\langle\beta^{-1}\rangle=\langle\beta\rangle$ which establishes symmetry. As for transitivity, if $\gamma\sim\delta$ and $\delta\sim\varepsilon$ then
\[\gamma\varepsilon^{-1}=\gamma(\delta^{-1}\delta)\varepsilon^{-1}=(\gamma\delta^{-1})(\delta\varepsilon^{-1})\in\langle\beta\rangle.\]
One can easily see that the equivalence class of an element $\gamma\in(\mathbb F_{q^2})^\times$ is $[\gamma]=\gamma\langle\beta\rangle=\{\gamma,\beta\gamma,\ldots,\beta^c\gamma\}$. Note further that $|[\gamma]|=c+1$. Indeed, if there are $0\leq j<i\leq c$ such that $\beta^i\gamma=\beta^j\gamma$ then $0=(\beta^i-\beta^j)\gamma=(\beta^{i-j}-1)\beta^j\gamma$, which is a contradiction as none of the factors in the product is zero.
In the following claim we show that, under some conditions on $\alpha,\beta$, the second property of axis-evasiveness is met by the construction above. We mention already here that the third condition in Definition35 is not met by the construction as is (regardless of the choice of $\alpha,\beta$), and we will alter it afterwards to meet that property as well.
\setcounter{nttheorem}{56}
\begin{ntclaim}[Author Claim57]
Assume that $\langle\beta\rangle\cap\ell_\alpha=\langle\beta\rangle\cap\ell_{\alpha^{-1}}=\varnothing$ and that $\langle\beta\rangle\cap\mathbb F_q=\{1\}$. Then, for every $\ell,\ell'\in\mathcal X$ (not necessarily distinct) it holds that $\ell'\cap(\ell)\leq c$.
\end{ntclaim}
\begin{proof}
Recall that $(\gamma)=\{\beta\gamma,\ldots,\beta^c\gamma\}$. Thus,
\[\bigcup_{\gamma\in\ell}(\gamma)=\bigcup_{\gamma\in\ell}\bigcup_{i=1}^c\{\beta^i\gamma\}=\bigcup_{i=1}^c\beta^i\ell.\]
Therefore,
\begin{equation}\ell'\cap(\ell)=\ell'\cap\bigcup_{\gamma\in\ell}(\gamma)=\bigcup_{i=1}^c(\ell'\cap\beta^i\ell).\tag{5.2}\end{equation}
Fix $i\in[c]$ and consider two cases. First, if $\ell$ is vertical, namely, $\ell=b+\alpha\mathbb F_q$ for some $b\in\mathbb F_q$, then $\beta^i\ell=\beta^ib+\alpha\beta^i\mathbb F_q$. Since, by assumption, $\langle\beta\rangle\cap\mathbb F_q=\{1\}$ we have that $\alpha\beta^i\mathbb F_q\ne\alpha\mathbb F_q$ and so the line $\beta^i\ell$ is not vertical. As, by assumption, $\langle\beta\rangle\cap\ell_{\alpha^{-1}}=\varnothing$, we have that $\alpha\beta^i\notin\mathbb F_q$ and so the line $\beta^i\ell$ is not horizontal either.
Second, consider the case that $\ell$ is horizontal $\ell=b\alpha+\mathbb F_q$ for some $b\in\mathbb F_q$. Then, $\beta^i\ell=b\alpha\beta^i+\beta^i\mathbb F_q$. Per our assumption that $\langle\beta\rangle\cap\ell_\alpha=\varnothing$, we have that $\beta^i\mathbb F_q\ne\alpha\mathbb F_q$ and so the line $\beta^i\ell$ is not vertical. As we assume $\langle\beta\rangle\cap\mathbb F_q=\{1\}$, we have that $\beta^i\mathbb F_q\ne\mathbb F_q$, and so the line $\beta^i\ell$ cannot be horizontal either. To summarize, we have that $\beta^i\ell\notin\mathcal X$. However, $\ell'\in\mathcal X$ and so $\beta^i\ell$ and $\ell'$ are two distinct lines. As such, the two lines intersect in at most one point. Equation(5.2) then yield $|\ell'\cap(\ell)|\leq c$.
\end{proof}
\paragraph{Informal discussion regarding the third property}
As mentioned above, the partition of $(\mathbb F_{q^2})^\times$ as defined above does not have the third property required for axis-evasiveness. Namely, there are $\gamma\in(\mathbb F_{q^2})^\times$ such that $[\gamma]$ intersects some axis-parallel line at more than one point. To get some idea on which equivalence classes $[\gamma]$ are problematic, let us first ask when do $\gamma,\beta\gamma$ are on some common axis-parallel line. We first observe that two points $\delta,\varepsilon\in(\mathbb F_{q^2})^\times$ are on a common axis-parallel line if and only if $\delta-\varepsilon\in\{1,\alpha\}\mathbb F_q$. Thus, $\gamma$ and $\beta\gamma$ are on the same axis-parallel line if and only if $\gamma-\beta\gamma=(1-\beta)\gamma\in\{1,\alpha\}\mathbb F_q$. This is equivalent to saying that $\gamma$ is on one of the two lines through the origin with slopes $1/(1-\beta),\alpha/(1-\beta)$.
More generally, $[\gamma]$ intersects with some axis-parallel line in more than one point if and only if $\beta^i\gamma-\beta^j\gamma\in\{1,\alpha\}\mathbb F_q$ for some $0\leq j<i\leq c$. Equivalently, $\gamma$ is on a line $\ell_\delta$ with
\begin{equation}\delta\in\left\{\frac1{\beta^i-\beta^j},\frac\alpha{\beta^i-\beta^j}\;\middle|\;0\leq j<i\leq c\right\}.\tag{5.3}\end{equation}
The key observation is that although there are a fair amount of ``bad'' points $\gamma$, they are all contained in a small number of lines. By ``small'' here we mean that the number is polynomial in $c$ and is independent of $q$. Thus, the hope is that by redefining the partition on these few problematic lines we will not harm the previous analysis by much. Indeed, no matter how we alter the partition restricted to these lines, if we make sure none of them is axis-parallel (by requiring more properties from $\alpha,\beta$) then each of these lines intersect an axis-parallel line at one point. As a result, the bound obtained in Claim57 will deteriorate proportionally to the number of lines above.
The only small technical issue is that even if $\gamma\in\ell_\delta$ for some slope $\delta$ as above, it is not the case that $[\gamma]\subseteq\bigcup_\varepsilon\ell_\varepsilon$ where $\varepsilon$ is taken from the set of slopes given by Equation(5.3). As we wish to alter the partition defined above, it would be cleaner to have all of the points in $[\gamma]$ of a problematic point $\gamma$ contained in the set of points on which we redefine the partition. Thus, we ``close'' the set of slopes given by Equation(5.3) to multiplication by $\beta$.
Ending the informal discussion and returning to the formal analysis, we consider the set of slopes.
\begin{equation}\Delta=\left\{\frac{\beta^k}{\beta^i-\beta^j},\frac{\alpha\beta^k}{\beta^i-\beta^j}\;\middle|\;0\leq j<i\leq c\text{ and }0\leq k\leq c\right\}.\tag{5.4}\end{equation}
Further define the set of all points in $(\mathbb F_{q^2})^\times$ covered by the lines with slopes from $\Delta$ by $U=\bigcup_{\delta\in\Delta}\ell_\delta$. This definition of $\Delta$ indeed fixes the technical caveat discussed above, as the following claim states.
\begin{ntclaim}[Author Claim58]
For every $\gamma\in(\mathbb F_{q^2})^\times$ either $[\gamma]\subseteq U$ or $[\gamma]\cap U=\varnothing$.
\end{ntclaim}
\begin{proof}
If an element $\varepsilon\in U$ then $\varepsilon\in\ell_\delta$ for some $\delta\in\Delta$. Note that $\beta\varepsilon\in\ell_{\beta\delta}$ and that $\beta\delta\in\Delta$. Hence, $\beta\varepsilon\in U$. Therefore, $\varepsilon\in U\Longrightarrow\varepsilon\langle\beta\rangle\subseteq U$. Assume now that $[\gamma]\cap U\ne\varnothing$, and take $\gamma\beta^i\in U$. By the above, $\gamma\beta^i\langle\beta\rangle\subseteq U$. The proof then follows as $\gamma\beta^i\langle\beta\rangle=\gamma\langle\beta\rangle=[\gamma]$.
\end{proof}
Define a new partition of $\mathbb F_{q^2}$ (including $0$) which agrees with the one that is given by Equation(5.1) on $\mathbb F_{q^2}^\times\setminus U$. By Claim58, this is well-defined. The new partition, restricted to $U$, is done as follows. Let $\delta_0\in\Delta$ be an arbitrary element. Note that
\[U=\ell_{\delta_0}\cup\bigcup_{\delta\in\Delta\setminus\{\delta_0\}}(\ell_\delta\setminus\{0\})\]
is a disjoint union. To partition $U$, we partition $\ell_{\delta_0}$ as well as each of $\ell_\delta\setminus\{0\}$ where $\delta\in\Delta\setminus\{\delta_0\}$ in an arbitrary way provided it has the least number of parts under the conditions that each part has size at most $c+1$. For ease of readability, we denote by $[\gamma]$ the class with respect to the new partition.
\begin{ntclaim}[Author Claim59]
Assume, on top of the assumptions of Claim57 that for every $\delta\in\Delta$, $\ell_\delta\notin\mathcal X$. Then, the new partition defined above is $(c,4c^2)$-axis evasive.
\end{ntclaim}
\begin{proof}
First, observe that by construction, every class intersects any axis-parallel line in at most one point. Indeed, classes that are outside of $U$ have this property by the definition of $U$ as can be easily verified (and discussed above). Moreover, by the way we redefined the partition restricted to $U$, every class that is a subset of $U$ is also a subset of a line $\ell_\delta$ for some $\delta\in\Delta$. As $\ell_\delta\notin\mathcal X$ by hypothesis, we have that the line and, as a result, the class it contains, intersects any axis-parallel line in at most one point. This establishes the third property of axis-evasiveness. The second property follows as, by construction, every part has size at most $c+1$.
Moving on to the second property, consider $\ell,\ell'\in\mathcal X$, not necessarily distinct. As outside of $U$ the partition is defined as before, Claim58 yields
\begin{equation}\left|\ell'\cap\bigcup_{\gamma\in\ell\setminus U}(\gamma)\right|\leq c.\tag{5.5}\end{equation}
Take $\gamma\in U\cap\ell$. Since, by construction $(\gamma)\subseteq\ell_\delta$ for some $\delta\in\Delta$, and since by hypothesis $\ell_\delta\notin\mathcal X$ we have that $|\ell'\cap\ell_\delta|=1$ and $(\gamma)\cap\ell'\subseteq\ell_\delta\cap\ell'$. Therefore, $|(\gamma)\cap\ell'|\leq1$. Together with Equation(5.5) we get that $|\ell'\cap(\ell)|\leq c+|U\cap\ell|$. Now, since $\ell\in\mathcal X$ and every line $\ell_\delta$ with slope $\delta\in\Delta$ is not in $\mathcal X$ we have that $|\ell\cap\ell_\delta|=1$. Thus, $|U\cap\ell|\leq|\Delta|$ which implies $|\ell'\cap(\ell)|\leq c+|\Delta|$.
To conclude the proof, we turn to bound $|\Delta|$. It is straightforward to give a bound of $O(c^3)$ though one can optimize the bound a bit. Indeed, with the notation of Equation(5.4), by multiplying by $\beta^{-\min(j,k)}$, one can rewrite
\begin{align}\Delta&=\left\{\frac1{\beta^i-\beta^j},\frac\alpha{\beta^i-\beta^j}\;\middle|\;0<j<i\leq c\right\}\notag\\&\quad\cup\left\{\frac{\beta^j}{\beta^i-1},\frac{\alpha\beta^j}{\beta^i-1}\;\middle|\;0<i\leq c,\ 0\leq j\leq c\right\}.\tag{5.6}\end{align}
Thus, $|\Delta|\leq3c^2$, and the proof follows.
\end{proof}
We summarize the discussion so far.
\begin{ntproposition}[Author Proposition60]
Let $\mathbb F_q$ be finite field. Let $h(x)\in\mathbb F_q[x]$ be a degree $2$ irreducible monic polynomial, and consider the field $\mathbb F_q[x]/\langle h(x)\rangle$ which we denote by $\mathbb F_{q^2}$. Let $\alpha,\beta\in\mathbb F_{q^2}$ be two elements satisfying:
\begin{enumerate}
\item $h(\alpha)=0$,
\item $\langle\beta\rangle\cap\mathbb F_q=\{1\}$,
\item $c+1=o(\beta)\leq\sqrt q/10$,
\item $\langle\beta\rangle\cap\ell_\alpha=\langle\beta\rangle\cap\ell_{\alpha^{-1}}=\varnothing$,
\item $(\langle\beta\rangle-\langle\beta\rangle)\cap\mathbb F_q=\{0\}$,
\item $(\langle\beta\rangle-\langle\beta\rangle)\cap\ell_\alpha=(\langle\beta\rangle-\langle\beta\rangle)\cap\ell_{\alpha^{-1}}=\{0\}$.
\end{enumerate}
Then, there exists a partition of $(\mathbb F_q)^2$ that is $(c,4c^2)$-axis-evasive, where $c=o(\beta)-1$. The number of parts in the partition is bounded above by $2q^2/(c+1)$.
\end{ntproposition}
To prove Proposition60 we need the following easy claim.
\begin{ntclaim}[Author Claim61]
Let $\delta\in(\mathbb F_{q^2})^\times$ be such that $\delta\notin\mathbb F_q\cup\ell_\alpha$ then, $\ell_\delta\notin\mathcal X$.
\end{ntclaim}
\begin{proof}
Write $\delta=a+\alpha b$ with $a,b\in\mathbb F_q$. Then, $\ell_\delta=(a+\alpha b)\mathbb F_q$. Observe that if $\ell_\delta$ is vertical then $a=0$ and so $\delta\in\ell_\alpha$. Similarly, if $\ell_\delta$ is horizontal then $b=0$ implying $\delta\in\mathbb F_q$.
\end{proof}
\begin{proof}[Author proof of Proposition60]
To bound the number of parts, recall that in the original partition, each part has size $c+1$. Moreover, in the altered partition we partition each line $\ell_\delta$ with slope $\delta\in\Delta$ (excluding the origin from all but for one of the lines $\ell_{\delta_0}$) to parts of size $c+1$ each, except for possibly one part. As $|\Delta|\leq3c^2$, the number of parts it bounded by
\[\frac{q^2-1}{c+1}+|\Delta|\left(1+\frac q{c+1}\right)\leq\frac{q^2-1}{c+1}+6cq\leq\frac{2q^2}c,\]
where the last inequality follows by our assumption that $o(\beta)\leq\sqrt q/10$.
To conclude the proof of the proposition, it suffices to show that for every $\delta\in\Delta$ it holds that $\ell_\delta\notin\mathcal X$. By Claim61, it suffices to prove that $\delta\notin\mathbb F_q\cup\ell_\alpha=\{1,\alpha\}\mathbb F_q$. There are two types of slopes $\delta\in\Delta$, according to whether they appear in the first or second set in Equation(5.6). The first kind is of the form
\[\delta=\frac{\alpha^k}{\beta^i-\beta^j},\]
with $0<j<i\leq c$ and $k\in\{0,1\}$. If $\delta\in\{1,\alpha\}\mathbb F_q$ then $\delta^{-1}\in\{1,\alpha^{-1}\}\mathbb F_q$ and so $\beta^i-\beta^j\in\{\alpha^k,\alpha^{k-1}\}\mathbb F_q$ in contradiction to our hypothesis. Consider now the other kind of slope
\[\delta=\frac{\alpha^k\beta^j}{\beta^i-1}\]
where $0<i\leq c$, $0\leq j\leq c$ and $k\in\{0,1\}$. If $\delta\in\{1,\alpha\}\mathbb F_q$ then $\delta^{-1}\in\{1,\alpha^{-1}\}\mathbb F_q$ and so $(\beta^i-1)\beta^{-j}\in\{\alpha^k,\alpha^{k-1}\}\mathbb F_q$. Note that $(\beta^i-1)\beta^{-j}=\beta^{i-j}-\beta^{-j}\in\langle\beta\rangle-\langle\beta\rangle$ and so we again get a contradiction.
\end{proof}
We are now ready to prove Theorem47. For the sake of readability, we repeat its statement here.
\begin{nttheorem}[Author Theorem62]
Let $P$ be a set of size $q$, where $q$ is an odd prime power. Let $c$ be an even integer such that $c+1\mid q+1$, and $c\leq\sqrt q/10$. Then, there exists a $(c,4c^2)$-axis evasive partition of $P^2$ with at most $2q^2/(c+1)$ parts.
\end{nttheorem}
\begin{proof}
As above, we identify $P^2$ with $\mathbb F_{q^2}$. It is a well-known fact that the multiplicative group $(\mathbb F_{q^2})^\times$ is cyclic. A basic result in group theory states that a cyclic group has a (unique) subgroup of every given size which divides the group size. Now, $|(\mathbb F_{q^2})^\times|=q^2-1=(q-1)(q+1)$. Thus, as $c+1\mid q+1$, there exists a subgroup $H$ of $(\mathbb F_{q^2})^\times$ of size $c+1$. The subgroup $H$ is cyclic, being a subgroup of a cyclic group. Let $\beta$ be a generator for $H$. We first prove that $\beta$ satisfies those hypothesis of Proposition60 that do not involve $\alpha$, namely, conditions(2) and(5).
\begin{ntclaim}[Author Claim63]
$(\langle\beta\rangle-\langle\beta\rangle)\cap\mathbb F_q=\{0\}$ and $\langle\beta\rangle\cap\mathbb F_q=\{1\}$.
\end{ntclaim}
\begin{proof}
Assume towards a contradiction that $\beta^i-\beta^j\in\mathbb F_q$ for some $0\leq j<i\leq c$. Since $x^q=x$ for every $x\in\mathbb F_q$, we get
\[\beta^i-\beta^j=(\beta^i-\beta^j)^q=\beta^{iq}-\beta^{jq},\]
where the last equality follows since $q$ is divisible by the characteristic of the field. Recall that $o(\beta)=c+1\mid q+1$ and so $\beta^{i(q+1)}=1$, implying $\beta^{iq}=\beta^{-i}$. Thus,
\[\beta^i-\beta^j=\frac1{\beta^i}-\frac1{\beta^j}=\frac{\beta^j-\beta^i}{\beta^{i+j}}.\]
As $\beta^i\ne\beta^j$ the above equation implies $\beta^{i+j}=-1$, and so $-1\in H$. Since $q$ is odd, the characteristic of the field $\mathbb F_{q^2}$ is odd and so $o(-1)=2$. Lagrange's Theorem then implies that $2\mid|H|=c+1$, which stands in contradiction to $c$ being even.
To prove that $\langle\beta\rangle\cap\mathbb F_q=\{1\}$, take $\beta^i$ with $0<i\leq c$. If $\beta^i\in\mathbb F_q$ then $\beta^{iq}=\beta^i$. On the other hand, we proved above that $\beta^{iq}=\beta^{-i}$, and so $\beta^i=\beta^{-i}$ implying $\beta^{2i}=1$. Therefore, $o(\beta)=c+1\mid2i$, but this is impossible as $0<i\leq c$ and, recall, $c$ is even.
\end{proof}
We proceed with the proof of Theorem47 by finding $\alpha\in\mathbb F_{q^2}$ that, together with the already chosen $\beta$, satisfies the remaining conditions in the hypothesis of Proposition60. Since $\mathbb F_{q^2}$ is a quadratic field extension of $\mathbb F_q$, every element $\gamma\in\mathbb F_{q^2}\setminus\mathbb F_q$ has degree $2$. That is, the minimal polynomial $h_\gamma$ of every such $\gamma$ over $\mathbb F_q$ is of degree $2$ (and can be made monic by dividing by the leading coefficient, if necessary). Indeed, $\deg(h_\gamma)$ cannot equal $1$ as this would imply $\gamma\in\mathbb F_q$. On the other hand,
\[2=[\mathbb F_{q^2}:\mathbb F_q]=[\mathbb F_{q^2}:\mathbb F_q(\gamma)][\mathbb F_q(\gamma):\mathbb F_q]=[\mathbb F_{q^2}:\mathbb F_q(\gamma)]\deg(h_\gamma),\]
which shows that if $\deg(h_\gamma)\ne1$ then $\deg(h_\gamma)=2$.
Thus, condition(1) in the hypothesis of Proposition60 holds for every element in $\mathbb F_{q^2}\setminus\mathbb F_q$. Hence, to prove that all the remaining conditions in the hypothesis of Proposition60 hold, it suffices to prove that there exists $\alpha\in\mathbb F_{q^2}\setminus\mathbb F_q$ which satisfies conditions(4) and(6). To this end, pick a set of ``slopes'' $\Delta'=\{\delta_1,\ldots,\delta_{q+1}\}\subseteq(\mathbb F_{q^2})^\times$ such that $(\mathbb F_{q^2})^\times$ is the disjoint union
\[(\mathbb F_{q^2})^\times=\bigcup_{\delta\in\Delta'}(\ell_\delta\setminus\{0\}).\]
For example, $\Delta'=\{a+\alpha\mid a\in\mathbb F_q\}\cup\{1\}$ will do. For $\delta\in(\mathbb F_{q^2})^\times$ let
\[I_\delta=|\langle\beta\rangle\cap\ell_\delta|+|(\langle\beta\rangle-\langle\beta\rangle)\cap(\ell_\delta\setminus\{0\})|.\]
Since the $\ell_\delta\setminus\{0\}$ with $\delta\in\Delta'$ are disjoint, $0\notin\langle\beta\rangle$, and since $|\langle\beta\rangle|=c+1$ and $|\langle\beta\rangle-\langle\beta\rangle|\leq(c+1)^2$, we have that
\[\mathbb E_\delta[I_\delta]\leq\frac{(c+1)^2+(c+1)}{q+1}\leq\frac{2(c+1)^2}{q+1},\]
where $\delta$ is sampled uniformly from $\Delta'$. By Markov's inequality, for at least $3/4$ of the elements $\delta\in\Delta'$ it holds that
\[|I_\delta|\leq\frac{8(c+1)^2}{q+1}.\]
Note that $(\mathbb F_{q^2})^\times$ is also a disjoint union of $\{\ell_{\delta^{-1}}\setminus\{0\}\mid\delta\in\Delta\}$. Thus, using the same argument as above, we get that for at least $1/2$ the elements $\delta\in\Delta'$, both $|I_\delta|$ and $|I_{\delta^{-1}}|$ are bounded by $8(c+1)^2/(q+1)$. But, as $c\leq\sqrt q/10$, this bound is strictly smaller than $1$, implying that $I_i=I_{q+1-i}=0$. That is, at least half the elements $\delta\in\Delta'$ satisfy conditions(4) and(6). Take $\alpha$ to be any of these elements. To conclude, we found $\alpha$ and $\beta$ for which all the conditions in the hypothesis of Proposition60 are met, and the proof follows.
\end{proof}

\begin{thebibliography}{43}
\bibitem{ref1}Noga Alon, Jeff Edmonds, and Michael Luby. Linear time erasure codes with nearly optimal recovery. In Proceedings of IEEE 36th Annual Foundations of Computer Science, pages 512–519. IEEE, 1995.
\bibitem{ref2}Noga Alon and Michael Luby. A linear time erasure-resilient code with nearly optimal recovery. IEEE Transactions on Information Theory, 42(6):1732–1736, 1996.
\bibitem{ref3}Sanjeev Arora, Carsten Lund, Rajeev Motwani, Madhu Sudan, and Mario Szegedy. Proof verification and the hardness of approximation problems. Journal of the ACM (JACM), 45(3):501–555, 1998.
\bibitem{ref4}Sanjeev Arora and Shmuel Safra. Probabilistic checking of proofs: A new characterization of NP. Journal of the ACM (JACM), 45(1):70–122, 1998.
\bibitem{ref5}László Babai, Lance Fortnow, Leonid A Levin, and Mario Szegedy. Checking computations in polylogarithmic time. In Proceedings of the twenty-third annual ACM symposium on Theory of computing, pages 21–32, 1991.
\bibitem{ref6}Laszlo Babai, Lance Fortnow, Noam Nisan, and Avi Wigderson. BPP has subexponential time simulations unless EXPTIME has publishable proofs. Computational Complexity, 3(4):307–318, 1993.
\bibitem{ref7}Omer Barkol, Yuval Ishai, and Ronny Roth. Locally decodable codes and their applications. PhD thesis, Computer Science Department, Technion, 2008.
\bibitem{ref8}Donald Beaver and Joan Feigenbaum. Hiding instances in multioracle queries. In Annual Symposium on Theoretical Aspects of Computer Science, pages 37–48. Springer, 1990.
\bibitem{ref9}Mihir Bellare and John Rompel. Randomness-efficient oblivious sampling. In Proceedings of the 35th Annual Symposium on Foundations of Computer Science, 1994, pages 276–287. IEEE, 1994.
\bibitem{ref10}Manuel Blum, Michael Luby, and Ronitt Rubinfeld. Self-testing/correcting with applications to numerical problems. In Proceedings of the twenty-second annual ACM symposium on Theory of computing, pages 73–83, 1990.
\bibitem{ref11}Benny Chor, Oded Goldreich, Eyal Kushilevitz, and Madhu Sudan. Private information retrieval. In Proceedings of IEEE 36th Annual Foundations of Computer Science, pages 41–50. IEEE, 1995.
\bibitem{ref12}Zeev Dvir. On matrix rigidity and locally self-correctable codes. computational complexity, 20(2):367–388, 2011.
\bibitem{ref13}Zeev Dvir, Parikshit Gopalan, and Sergey Yekhanin. Matching vector codes. SIAM Journal on Computing, 40(4):1154–1178, 2011.
\bibitem{ref14}Klim Efremenko. 3-query locally decodable codes of subexponential length. SIAM Journal on Computing, 41(6):1694–1703, 2012.
\bibitem{ref15}Peter Gemmell, Richard Lipton, Ronitt Rubinfeld, Madhu Sudan, and Avi Wigderson. Self- testing/correcting for polynomials and for approximate functions. In STOC, volume 91, pages 32–42. Citeseer, 1991.
\bibitem{ref16}Peter Gemmell and Madhu Sudan. Highly resilient correctors for polynomials. Information processing letters, 43(4):169–174, 1992.
\bibitem{ref17}Oded Goldreich. A sample of samplers: A computational perspective on sampling. In Studies in Complexity and Cryptography, pages 302–332. Springer, 2011. doi:10.1007/ 978-3-642-22670-0\_24.
\bibitem{ref18}Oded Goldreich and Madhu Sudan. Locally testable codes and PCPs of almost-linear length. Journal of the ACM (JACM), 53(4):558–655, 2006.
\bibitem{ref19}Sivakanth Gopi, Swastik Kopparty, Rafael Oliveira, Noga Ron-Zewi, and Shubhangi Saraf. Locally testable and locally correctable codes approaching the gilbert-varshamov bound. IEEE Transactions on Information Theory, 64(8):5813–5831, 2018.
\bibitem{ref20}Alan Guo, Swastik Kopparty, and Madhu Sudan. New affine-invariant codes from lifting. In Proceedings of the 4th conference on Innovations in Theoretical Computer Science, pages 529–540. ACM, 2013.
\bibitem{ref21}Venkatesan Guruswami, Atri Rudra, and Madhu Sudan. Essential coding theory. Draft available at http://www.cse.buffalo.edu/\textasciitilde{}atri/courses/coding-theory/book, 2012.
\bibitem{ref22}Brett Hemenway, Rafail Ostrovsky, and Mary Wootters. Local correctability of expander codes. Information and Computation, 243:178–190, 2015.
\bibitem{ref23}Jonathan Katz and Luca Trevisan. On the efficiency of local decoding procedures for error- correcting codes. In Proceedings of the thirty-second annual ACM symposium on Theory of computing, pages 80–86, 2000.
\bibitem{ref24}Tali Kaufman and Madhu Sudan. Sparse random linear codes are locally decodable and testable. In 48th Annual IEEE Symposium on Foundations of Computer Science (FOCS’07), pages 590–600. IEEE, 2007.
\bibitem{ref25}Kiran S Kedlaya and Sergey Yekhanin. Locally decodable codes from nice subsets of finite fields and prime factors of mersenne numbers. SIAM Journal on Computing, 38(5):1952–1969, 2009.
\bibitem{ref26}Swastik Kopparty, Or Meir, Noga Ron-Zewi, and Shubhangi Saraf. High-rate locally correctable and locally testable codes with sub-polynomial query complexity. Journal of the ACM (JACM), 64(2):11, 2017.
\bibitem{ref27}Swastik Kopparty, Shubhangi Saraf, and Sergey Yekhanin. High-rate codes with sublinear-time decoding. Journal of the ACM (JACM), 61(5):28, 2014.
\bibitem{ref28}Ray Li and Mary Wootters. Lifted multiplicity codes and the disjoint repair group property. In Approximation, Randomization, and Combinatorial Optimization. Algorithms and Techniques (APPROX/RANDOM 2019). Schloss Dagstuhl-Leibniz-Zentrum fuer Informatik, 2019.
\bibitem{ref29}Richard J. Lipton. Efficient checking of computations. In Annual Symposium on Theoretical Aspects of Computer Science, pages 207–215. Springer, 1990.
\bibitem{ref30}Irving S Reed. A class of multiple-error-correcting codes and the decoding scheme. Technical report, Massachusetts inst of tech Lexington Lincoln lab, 1953.
\bibitem{ref31}Omer Reingold, Salil Vadhan, and Avi Wigderson. Entropy waves, the zig-zag graph product, and new constant-degree expanders and extractors. In Electronic Colloquium on Computational Complexity (ECCC), page 18, 2001. URL: https://eccc.weizmann.ac.il/report/2001/018/.
\bibitem{ref32}Ronitt Rubinfeld and Madhu Sudan. Robust characterizations of polynomials with applications to program testing. SIAM Journal on Computing, 25(2):252–271, 1996.
\bibitem{ref33}C. E. Shannon. A mathematical theory of communication. ACM SIGMOBILE Mobile Computing and Communications Review, 5(1):3–55, 2001. Originally appeared in Bell System Tech. J. 27:379–423, 623–656, 1948.
\bibitem{ref34}Carl Siegel. Über die classenzahl quadratischer zahlkörper. Acta Arithmetica, 1(1):83–86, 1935.
\bibitem{ref35}Madhu Sudan, Luca Trevisan, and Salil Vadhan. Pseudorandom generators without the XOR lemma. Journal of Computer and System Sciences, 62(2):236–266, 2001.
\bibitem{ref36}Luca Trevisan. List-decoding using the XOR lemma. In 44th Annual IEEE Symposium on Foundations of Computer Science, 2003. Proceedings., pages 126–135. IEEE, 2003.
\bibitem{ref37}Arnold Walfisz. Zur additiven zahlentheorie. ii. Mathematische Zeitschrift, 40(1):592–607, 1936.
\bibitem{ref38}David Woodruff. New lower bounds for general locally decodable codes. In Electronic Colloquium on Computational Complexity (ECCC), volume 14, 2007.
\bibitem{ref39}S. Yekhanin. Locally decodable codes. In International Computer Science Symposium in Russia, pages 289–290. Springer, 2011.
\bibitem{ref40}Sergey Yekhanin. Towards 3-query locally decodable codes of subexponential length. Journal of the ACM (JACM), 55(1):1–16, 2008.
\bibitem{ref41}Kalina Petrova Zeev Dvir. Lecture 1: Introduction. Lecture notes: https://www.cs. princeton.edu/\textasciitilde{}zdvir/LDCnotes/LDC1.pdf, year = 2016.
\bibitem{ref42}Kalina Petrova Zeev Dvir. Lecture 4: Lower bounds for r-query LDCs. Lecture notes: https://www.cs.princeton.edu/\textasciitilde{}zdvir/LDCnotes/LDC4.pdf, 2016.
\bibitem{ref43}Victor Vasilievich Zyablov. An estimate of the complexity of constructing binary linear cascade codes. Problemy Peredachi Informatsii, 7(1):5–13, 1971.
\end{thebibliography}
\end{document}
